% !TeX program = xelatex
% ============================================================================
% 2026 高教社杯 C 题 —— 微网购电与储能调度 论文 v2
%
% 手动构建（在 paper/ 目录下执行）：
%   xelatex main.tex && xelatex main.tex
% 连续编译两次（交叉引用需要第二遍）。产物为 main.pdf。
% 依赖：TeX Live / MacTeX（ctex、pgfplots、listings、booktabs 等均为标准组件）。
%
% 合规提示（2026 届格式规范）：
%   - 电子版第一页必须为摘要专用页：本文件已按此排版，页码自摘要页起、页脚中部。
%   - 纸质版需在摘要页前另加官方承诺书与编号专用页（不写入本文件）。
%   - 摘要页/正文/附录不得出现任何身份、学校、赛区信息。
%
% 图件：正文图取 results/paper/（fig01–17，paperfigs 命令重绘），题 4 两图为
% pgfplots 矢量绘制；附录图 figA2。附录 B 为外生参数标定，附录 C 收 5 个核心建模文件。
% ============================================================================
\documentclass[12pt,a4paper]{ctexart}
\usepackage[margin=2.5cm]{geometry}
\usepackage{amsmath,amssymb,bm}
\usepackage{graphicx}
\usepackage{booktabs}
\usepackage{multirow}
\usepackage{float}
\usepackage{caption}
\usepackage{xcolor}
\usepackage{listings}
\usepackage{enumitem}
\usepackage{pgfplots}
\pgfplotsset{compat=1.16}
\usetikzlibrary{arrows.meta,positioning}
\graphicspath{{../results/paper/}{../results/}}

\captionsetup{font=small}
\setlength{\parskip}{1pt}
% 2026 规范：页码自摘要页起、页脚中部；全文不用页眉
\pagestyle{plain}

% 中文章节号（一、二、三……），小节用阿拉伯数字
\renewcommand{\thesection}{\chinese{section}}
\renewcommand{\thesubsection}{\arabic{section}.\arabic{subsection}}

% 代码样式
\lstset{
  language=Python,
  basicstyle=\ttfamily\zihao{-5},
  keywordstyle=\color{blue!70!black},
  commentstyle=\color{gray},
  stringstyle=\color{green!40!black},
  breaklines=true,
  frame=single,
  numbers=left,
  numberstyle=\tiny\color{gray},
  showstringspaces=false,
  tabsize=4,
  columns=flexible,
}

% 代码附录中的 Unicode 符号映射（literate 在 XeLaTeX 下对多字节字符不生效，
% 用 newunicodechar 把字符变为 active，可在 listings 中正常展开）
\usepackage{newunicodechar}
\newunicodechar{≤}{\ensuremath{\leq}}
\newunicodechar{≥}{\ensuremath{\geq}}
\newunicodechar{∧}{\ensuremath{\wedge}}
\newunicodechar{∨}{\ensuremath{\vee}}
\newunicodechar{⁺}{\ensuremath{^{+}}}
\newunicodechar{⁻}{\ensuremath{^{-}}}
\newunicodechar{−}{\ensuremath{-}}
\newunicodechar{·}{\ensuremath{\cdot}}
\newunicodechar{≈}{\ensuremath{\approx}}
\newunicodechar{→}{\ensuremath{\rightarrow}}
\newunicodechar{×}{\ensuremath{\times}}
\newunicodechar{①}{\textcircled{\scriptsize 1}}
\newunicodechar{②}{\textcircled{\scriptsize 2}}
\newunicodechar{③}{\textcircled{\scriptsize 3}}
\newunicodechar{④}{\textcircled{\scriptsize 4}}
\newunicodechar{⑤}{\textcircled{\scriptsize 5}}
\newunicodechar{⑥}{\textcircled{\scriptsize 6}}
% 希腊字母与数学符号（代码注释中出现）
\newunicodechar{α}{\ensuremath{\alpha}}
\newunicodechar{β}{\ensuremath{\beta}}
\newunicodechar{γ}{\ensuremath{\gamma}}
\newunicodechar{δ}{\ensuremath{\delta}}
\newunicodechar{ε}{\ensuremath{\varepsilon}}
\newunicodechar{θ}{\ensuremath{\theta}}
\newunicodechar{κ}{\ensuremath{\kappa}}
\newunicodechar{λ}{\ensuremath{\lambda}}
\newunicodechar{ν}{\ensuremath{\nu}}
\newunicodechar{π}{\ensuremath{\pi}}
\newunicodechar{ρ}{\ensuremath{\rho}}
\newunicodechar{σ}{\ensuremath{\sigma}}
\newunicodechar{τ}{\ensuremath{\tau}}
\newunicodechar{φ}{\ensuremath{\varphi}}
\newunicodechar{η}{\ensuremath{\eta}}
\newunicodechar{ℓ}{\ensuremath{\ell}}
\newunicodechar{∈}{\ensuremath{\in}}
\newunicodechar{≡}{\ensuremath{\equiv}}
\newunicodechar{≠}{\ensuremath{\neq}}
\newunicodechar{⇒}{\ensuremath{\Rightarrow}}
\newunicodechar{′}{\ensuremath{'}}
\newunicodechar{⚠}{\textbf{!}}
\newunicodechar{₁}{\ensuremath{_{1}}}
\newunicodechar{⁰}{\ensuremath{^{0}}}
\newunicodechar{²}{\ensuremath{^{2}}}
\newunicodechar{³}{\ensuremath{^{3}}}
\newunicodechar{⁴}{\ensuremath{^{4}}}
\newunicodechar{⁵}{\ensuremath{^{5}}}
\newunicodechar{⁶}{\ensuremath{^{6}}}
\newunicodechar{️}{}

% 常用记号
\newcommand{\Zon}{Z^{\mathrm{online}}}
\newcommand{\ZCE}{Z^{\mathrm{CE}}}
\newcommand{\Zep}{Z^{\mathrm{ex\text{-}post}}}
\newcommand{\ZPI}{Z^{\mathrm{PI}}}
\newcommand{\ZD}{Z^{\mathrm{D}}}
\newcommand{\Zmis}{Z^{\mathrm{mis}}}
\newcommand{\nuend}{\nu^{\mathrm{end}}}

\begin{document}

% ===================== 摘要专用页（电子版第一页，页码自此起） =====================
\thispagestyle{plain}
\vspace*{-2em}
\begin{center}
{\Large\heiti 基于两阶段随机规划与在线随机模型预测控制的\\[4pt]微网购电与储能调度模型}
\end{center}
\vspace{1.5em}
\begin{center}\large\heiti 摘\quad 要\end{center}

针对含光伏与储能微网的购电费用最小化问题，本文以“供需平衡方程 + 储能荷电状态（SOC）递推”为四问共用的物理内核，按信息结构与罚则的逐级收紧，建立“预报—规划—执行—结算”四层统一优化框架，并用 2025 全年 334 个交付日做无信息泄漏的滚动回测。

针对问题一（典型日、完美信息），建立以购电费最小为目标的线性规划，最优购电费为 \textbf{35\,118.60 元}，较无电池直接购电节省 \textbf{26.92\%}，较阈值规则策略节省 \textbf{4.34\%}，昼夜循环水位 8\,550 kWh。由互补松弛条件推出套利价格门槛 $1/\eta^2\approx1.23$（实际峰谷比 3.76 远超之），并证明充放电天然互斥，模型无需 0--1 变量即可保持纯线性规划。

针对问题二（负载光伏逐日变化、断供按 5 倍电价紧急购电），认定 0:00 仅掌握历史信息，建立两阶段场景线性规划产生计划购电量、在线随机模型预测控制（MPC）逐时段执行。全年交付费用 \textbf{13\,745\,816 元}，紧急购电占 6.74\%，缺口频率 4.30\% 与报童临界频率 4.5\% 一致；同口径下确定性等价 MPC 年贵 \textbf{50.9 万元}，说明未来必须以概率分布而非单条路径进入模型；无电池对照贵 24.4\%，题一计划复读贵 80.5\%。

针对问题三（一天四次光伏预报、偏差按 50\%/150\% 结算），建立分块承诺制修订模型：0:00 报计划、6/12/18 时逐块锁定调整量，偏差费化简为 $\tfrac12 p_t|a_t-g_t|$；光伏预报采用反方差融合，误差较官方预报直用下降 36.6\%。全年费用 \textbf{13\,023\,147 元}，较问题二再降 \textbf{5.25\%}，紧急费占比压至 0.91\%；修订权逐级增量 $-28.8/-16.9/-15.3$ 万元、合计 \textbf{61.0 万元（4.48\%）}，其中 6:00 档贡献最大。

针对问题四（实时波动电价），认定附件四为实现值（0:00 需自建价格预报），建立负载—光伏—电价三元联合情景模型，末端估值逐日标定（零超参数）。全年回测：计划线费用 \textbf{14\,447\,040 元}、调整线 \textbf{13\,671\,435 元}（再省 \textbf{5.37\%}）；电价预报 MAPE 7.51\%，“夜间为当日最便宜时段块”预报 365/365 命中；日末水位呈“次日周五 2\,594 对其他 6\,974 kWh”的星期台阶，验证逐日化的必要性；常数价退化检验下两条交付线零误差复现问题二与问题三结果。

模型检验方面：对账恒等式与能量守恒残差均不超过 $10^{-6}$，序关系全年零违例，报童日频超买占比 88.6\%/86.2\% 落入理论区间 $80\%\pm10\%$；末端估值扰动 $0.60\sim1.00$ 元/kWh 时费用仅升 $0.79\%\sim1.43\%$。围绕主模型另建一套“只改一处”的副模型体系，分别量化电池（$+24.4\%$）、优化（$+16.5\%$）、逐日适应（$+80.5\%$）、执行层机制（年 50.9 万元）与修订权（$4.48\%$）的价值。

\vspace{0.5em}
\noindent\textbf{关键词：}微网；储能调度；两阶段随机规划；随机模型预测控制；反方差融合；对照实验

\newpage

% ===================== 一、问题重述 =====================
\section{问题重述}

在“双碳”目标与新型电力系统建设背景下，光伏等分布式能源快速发展，但其强间歇性使大规模直接并网冲击电网安全。微网作为解决分布式能源本地消纳的小型发用电系统，可以在外网电价较低或光伏有盈余的时段为储能充电、在电价较高的时段放电，利用峰谷价差降低购电成本。某居民小区自建微网：光伏阵列与一块储能电池同接于内部母线，并可向外网购电；系统拓扑与能量流转如图 \ref{fig:bus} 所示。运营目标是在保证小区负载不断电的前提下，最小化全年购电费用。储能参数由附录一给定：标称容量 12\,000 kWh，运行区间 1\,200--10\,800 kWh，最大充放电功率 5\,000 kW，充放电单向效率均为 0.90。题目给出四个逐级加难的问题：

\begin{figure}[H]
\centering
\begin{tikzpicture}[scale=0.76, every node/.style={transform shape},
    box/.style={draw, rectangle, minimum width=2.2cm, minimum height=1cm, align=center, thick},
    busbox/.style={draw, rectangle, minimum width=2.5cm, minimum height=1.6cm, align=center, thick},
    arrow/.style={-{Stealth[length=3mm, width=2mm]}, thick},
    dashedarrow/.style={-{Stealth[length=3mm, width=2mm]}, thick, dashed},
    labeltext/.style={font=\small\itshape}
]
    \node[busbox] (bus) {母线};
    \node[box, left=2cm of bus] (grid) {外网\\(购电付费)};
    \node[box, below=1.2cm of grid] (pv) {光伏\\(免费注入)};
    \node[box, right=3cm of bus] (load) {小区负载};
    \node[box, below=1.2cm of load] (curt) {弃电\\(免费耗散)};
    \node[box, below=1.5cm of bus] (inv) {逆变器\\[2pt] \(\eta = 0.9\)};
    \node[box, below=1.5cm of inv] (bat) {电池\\[2pt] \(S_t\)};
    \draw[arrow] (grid) -- node[midway, above, labeltext] {\(g_t\)} (bus);
    \draw[arrow] (pv) -- node[midway, above, labeltext] {\(r_t\)} (bus);
    \draw[arrow] (bus) -- node[midway, above, labeltext] {\(\ell_t\)} (load);
    \draw[dashedarrow] (bus) -- node[midway, above, labeltext] {\(w_t\)} (curt);
    \draw[arrow] ([xshift=-4pt]bus.south) -- node[midway, left, labeltext] {\(c_t\)} ([xshift=-4pt]inv.north);
    \draw[arrow] ([xshift=4pt]inv.north) -- node[midway, right, labeltext] {\(d_t\)} ([xshift=4pt]bus.south);
    \draw[arrow] ([xshift=-4pt]inv.south) -- node[midway, left, labeltext] {\(\eta c_t\)} ([xshift=-4pt]bat.north);
    \draw[arrow] ([xshift=4pt]bat.north) -- node[midway, right, labeltext] {\(d_t/\eta\)} ([xshift=4pt]inv.south);
    \node[right=1cm of load, align=left, font=\footnotesize] (notes) {
        \textbf{母线侧计量：}\\
        单时段充放电量\\ \(\le 833.33\) kWh\\[6pt]
        \textbf{效率损耗发生处：}\\ \(\eta = 0.9\)\\[6pt]
        \textbf{电池侧计量：}\\ \(\eta c_t\) 落袋，\(d_t/\eta\) 取出\\[6pt]
        \textbf{电池运行区间：}\\ \(1200 \le S_t \le 10800\) kWh
    };
    \draw[dashed, gray] ([xshift=-1.5cm, yshift=0.5cm]inv.north west) rectangle ([xshift=1.5cm, yshift=-0.5cm]bat.south east);
\end{tikzpicture}
\caption{微网系统物理拓扑与电量流转图}
\label{fig:bus}
\end{figure}


\begin{enumerate}[leftmargin=2em]
  \item \textbf{问题一（确定性日前调度）}：电价与负载每天相同，已知光伏出力一日预测；储能 0:00 与 24:00 储电量相等。在每天 0:00 制定当日购电与充放电计划。
  \item \textbf{问题二（波动负载与紧急购电）}：负载与光伏随日期变化（附件二给出全年实际值），供电不足时按交易时刻电价 5 倍紧急购电。逐日制定 2 月 1 日至 12 月 31 日的购电策略。
  \item \textbf{问题三（多时刻预报与偏差结算）}：每天 0:00、6:00、12:00、18:00 可获得未来 24 小时整点光伏预报（附件三）；计划高于调整的部分按 50\% 电价结算违约，调整高于计划的部分按 150\% 加价。制定计及预报修订的策略，并分析是否需要引入其他时刻的预报。
  \item \textbf{问题四（实时波动电价）}：外网电价实时波动（附件四给出全年逐时段实现值），在波动电价下重新求解问题二与问题三。
\end{enumerate}

各问分别交付 result1--result4 共五个结果文件，报表粒度为 10 分钟（购电）与 4 小时（充放电）。

% ===================== 二、问题分析 =====================
\section{问题分析}

\subsection{总体结构：一条物理内核，四级信息收紧}

四问共享同一物理系统，差别只在“决策时知道什么、做错了罚多少”。因此本文不分别为四问各建一个模型，而是抽取一个共用的物理可行域（母线能量平衡 + SOC 递推 + 功率与容量边界），在其上逐级改变信息结构与结算规则（图 \ref{fig:framework}）：问题一是完美信息下的单阶段线性规划；问题二引入两阶段决策（0:00 锁定计划、日内实时补救）与 5 倍断供罚则，对应两阶段随机规划\cite{birge}；问题三把“一次锁定”放宽为分块承诺下的嵌套滚动时域优化——外环在 6/12/18 时逐块锁定调整量、内环逐时段滚动执行，偏差罚则对称；问题四把电价由确定参数变为第三个不确定源。这样的结构使前一问的解是后一问的严格特例（如问题三中令 $a\equiv g$ 即退化为问题二），从而“后一问的费用不得超过前一问”本身成为判定实现正确性的检验。主模型给出交付费用之后，第 \ref{sec:aux} 节用同一套执行层构造只改一处结构的副模型，把“电池、优化、信息、机制各值多少钱”从叙事变成可复核的差额。

\subsection{问题一分析}

问题一的数学本质是单阶段线性优化。容易想到引入 0--1 变量刻画“同一时段不能既充又放”的物理常识，使模型成为混合整数规划；但本文证明（第 \ref{sec:q1} 节性质 1）在最优解处充放天然互斥，0--1 变量可以整组移除，模型保持纯线性规划。这不只是简化：全年 334 天滚动求解（问题二至四）的计算可行性直接建立在该结论上。同时，由对偶变量可推出储能套利是否划算的价格门槛 $1/\eta^2$，为解释解的结构提供理论标尺\cite{jiang2018,sikou,sioshansi}。

\subsection{问题二分析}

问题二的核心困难是信息：题面要求 0:00 报出全天计划，却没有说明当时掌握哪些信息。若假定已知当日真值，5 倍罚则永远不会触发，模型退化为问题一，与题意不符。本文认定 0:00 仅掌握截至昨日的全部历史实测（附件二按日解锁），全年做诚实回测。在此信息结构下，计划购电量是“看不到未来”的第一阶段决策，天然对应两阶段随机规划的样本平均近似（SAA）\cite{birge}；执行层若把未来压缩成一条确定性参考路径（确定性等价），会把“可能缺口”当作“必然缺口”处理，本文的回测显示该做法年费用高出 50.9 万元，因此执行层采用保留情景分布的在线随机 MPC\cite{rawlings}。情景覆盖水平可用报童模型的临界分位解释\cite{qin2011}。

\subsection{问题三分析}

问题三的表象是“多了三次改单机会”，实质是结算规则与信息结构同时变化：计划量 $g$ 自 0:00 起降为计费基准，实际注入是调整量 $a$。若允许任意时刻反复修订，决策维度为 $4\times144$；本文证明每个时段只在末次修订机会调整一次即弱占优（块分解定理），把问题约简为四个 36 维块决策。预报侧，官方预报误差随提前期急剧增大，而历史持续法误差与提前期无关，两者互补，这提示采用最小方差意义下的反方差加权融合\cite{bates1969}，而非简单选用其一。

\subsection{问题四分析}

问题四的关键认定是附件四的地位：它与附件二句式同构、同为全年实现值，故 0:00 同样不可知当日电价，需自建价格预报（若允许 0:00 读取当日 144 个实现价，则同一论文对两份同构附件采用了两套信息纪律）。该读法可行的依据有二：其一，电价乘性分解（日水平 $\times$ 日内形态）解释对数方差的 92.32\%，日内形态高度稳定；其二，可证明日内调度只依赖价格形态与跨日价格比，与日价格水平无关。价格与负载残差相关系数达 0.89，最坏组合（高缺口 $\times$ 高价）必须由三元联合情景刻画，独立抽样会系统性低估紧急购电费用。

\begin{figure}[H]
\centering
\includegraphics[width=0.80\textwidth]{fig02_三问递进-信息决策罚则.png}
\caption{四问递进关系：信息结构、决策结构与罚则逐级收紧，前一问是后一问的严格特例}
\label{fig:framework}
\end{figure}

% ===================== 三、模型假设与符号说明 =====================
\section{模型假设与符号说明}

\subsection{符号说明}

\begin{table}[H]
\centering\footnotesize
\caption{主要符号（左：集合、数据与常数；右：决策变量与口径）}
\begin{tabular}{clcl}
\toprule
符号 & 含义 & 单位 & 备注 \\
\midrule
\multicolumn{4}{c}{\textbf{集合与索引}} \\
$t,\ \mathcal{T}$ & 日内 10 分钟时段及集合 & -- & $|\mathcal{T}|=T=144$ \\
$j$ & 日索引 & -- & 报告期共 334 天 \\
$m,\ \tau_m,\ \mathcal{B}_m$ & 题\,三修订阶段 / 时刻 / 块 & -- & 块长 36 时段 \\
\midrule
\multicolumn{4}{c}{\textbf{已知数据与常数}} \\
$p_t$ & 外网购电单价 & 元/kWh & 题四为波动实现价 $p_t^{(j)}$ \\
$\ell_t,\ r_t$ & 负载需电量 / 光伏发电量 & kWh & 功率 $\times\Delta$ 预折算 \\
$n_t$ & 净负荷 $\ell_t-r_t$ & kWh & \\
$\Delta$ & 单时段时长 & h & $1/6$ \\
$\eta$ & 单向充放效率 & -- & 0.90，往返 $\eta^2=0.81$ \\
$\kappa$ & 紧急购电倍率 & -- & 5 \\
$\beta^-,\ \beta^+$ & 题\,三违约 / 加价倍率 & -- & 0.5 / 1.5 \\
\bottomrule
\end{tabular}\hfill
\begin{tabular}{clcl}
\toprule
符号 & 含义 & 单位 & 备注 \\
\midrule
\multicolumn{4}{c}{\textbf{决策变量}} \\
$g_t$ & 计划购电量 & kWh & 0:00 锁定 \\
$a_t$ & 调整购电量（题三） & kWh & 实际注入量 \\
$c_t,\ d_t$ & 充电 / 放电量（母线侧） & kWh & $\le 833.33$ \\
$u_t$ & 紧急购电量 & kWh & 单价 $\kappa p_t$ \\
$w_t$ & 弃电量（平衡松弛） & kWh & 题四拆弃光与超买 \\
\midrule
\multicolumn{4}{c}{\textbf{状态、参数与口径}} \\
$S_t$ & 时段末储电量 & kWh & $1200\le S_t\le 10800$ \\
$\nuend$ & 日末储电量估值 & 元/kWh & 题\,四逐日标定 \\
$\varphi(h)$ & 反方差融合权重（题三） & -- & 提前期 $h$ 的函数 \\
$\alpha^\star$ & 报童临界分位 & -- & $(\kappa-1)/(\kappa-\theta)$ \\
$\Zon$ 等 & 全年费用口径 & 元 & online / CE / 事后 / PI \\
\bottomrule
\end{tabular}
\label{tab:sym}
\end{table}

\subsection{模型假设}

假设按用途分三类，每条均在后文被使用。

\textbf{物理与口径类}
\begin{enumerate}[label=A\arabic*., leftmargin=3em]
  \item 光伏盈余弃电允许且免费，弃电量 $w_t$ 不设上界。该假设使“供电不低于负载”的不等式约束可以写成等式（$w_t$ 为松弛量），且是充放互斥性证明的承重条件；
  \item 微网不可向大电网售电，$g_t\ge0$ 且无售电变量（题面无上网机制）；
  \item 外网购电无功率上限（题面未给限制）；储能功率上限 5\,000 kW 与容量区间按附录一执行；
  \item 模型内部充放变量按母线侧计量（与平衡方程、功率上限同口径），报表按电池侧计量（$\eta c_t$、$d_t/\eta$，与储电量对账自洽），换算仅在导出报表时发生一次。
\end{enumerate}

\textbf{信息结构类}
\begin{enumerate}[label=A\arabic*., leftmargin=3em, start=5]
  \item \textbf{诚实回测}：第 $j$ 日 0:00 仅掌握截至第 $j-1$ 日的全部历史实测、附件一与日历；附件二（负载光伏）与附件四（电价）均为按日解锁的实现值，当日真值不可见。这是全题最关键的认定，理由：附件二与附件四在附录二中句式同构，若对前者不许偷看、对后者允许，则信息纪律自相矛盾；
  \item 两阶段决策结构：计划购电量 $g_t$ 于 0:00 锁定、全额计费、不可撤改（题三起可在规定时刻按偏差价调整）；充、放、紧急购电、弃电为实时补救量。依据是题面对三张报表的措辞差异（“计划购电量”冠“0:00 制定”，充放电与紧急购电未冠）。
\end{enumerate}

\textbf{规则与参数类}
\begin{enumerate}[label=A\arabic*., leftmargin=3em, start=7]
  \item 紧急购电是断供兜底：无功率上限、单价 $\kappa p_t$，但不得超过当步计划缺口，即 $u_t\le(n_t-g_t)^+$（题三为 $n_t-a_t$）。该约束不含任何待标定参数，可证等价于“紧急购电时不得充电”；无此约束时，5 倍价电会充当日内采购通道（见 \ref{sec:q2cap} 节）；
  \item 题\,三结算按“最终调整量对 0:00 计划量”单次比对计偏差费，罚率挂该电量所属时段电价（题面“交易时刻电价”）；块分解定理保证该读法对链式读法不敏感；
  \item 题\,四价格预报自建（近 4 个同星期几同时段均值），不做日内更新与偏差校正（实测改善不足 1.5 个百分点，不抵一条代码路径）；
  \item 一月为暖机期：用同一模型完整滚动但不计入报告（题面初始电量为 1 月 1 日，交付自 2 月 1 日起，一月空档由此补齐）；全部随机性固定种子。
\end{enumerate}

% ===================== 四、数据处理与核验 =====================
\section{数据处理与核验}

\subsection{量纲与对齐}

输入功率（kW）统一乘 $\Delta=1/6$ 折算为电量（kWh），功率上限 5\,000 kW 折算为单时段 833.33 kWh，模型内部不再出现时间步长。附件行标签为时段末点，与 result 模板整体错一格，按行号对位填写（模板标签不改动）。四个数据附件均无缺失、无非正值，不做填补与剔除；题四最低价 0.0076 元/kWh 落在午间光伏出力高峰（光伏挤出效应的机理见文献\cite{denholm2007}），视为信号而非噪声。

\subsection{输入数据的形态}

附件一典型日呈“双谷双峰”电价与昼间光伏结构（图 \ref{fig:data}）：电价谷段（00:00--06:00、11:00--13:00）低至约 0.42 元/kWh，峰段（06:00--10:00、17:00--21:00）高达 1.40 元/kWh，峰谷比 3.76；光伏出力集中于 04:40--19:20，日发电 55\,483 kWh，恰为日负荷 111\,025 kWh 的一半；10:00--14:00 光伏大于负载出现盈余，为储能充电提供天然空间。全年尺度上（图 \ref{fig:year}），负载与光伏呈夏丰冬欠的季节结构，负载有清晰的周内节律（周五、周六约为其余五天的 63\%）——该节律是后文负载预报与电价预报的主要信息来源。

\begin{figure}[H]
\centering
\includegraphics[width=0.80\textwidth]{fig03_附件1典型日-电价负荷光伏.png}
\caption{附件一典型日：电价双谷双峰、光伏昼间出力与净负荷盈余区间}
\label{fig:data}
\end{figure}

\begin{figure}[H]
\centering
\includegraphics[width=0.80\textwidth]{fig08_附件2全年负荷光伏.png}
\caption{附件二全年负载与光伏：季节结构与周内节律（暖机期一月不计报告）}
\label{fig:year}
\end{figure}

\subsection{附件三的对齐与精度}

“预报 $h$ 小时”认定为对目标整点的瞬时功率。四种对齐假设中以 RMSE 判别：采纳方案 370.6 kW（相关系数 0.9921）明显优于次优（418.4 kW）。整点预报以已实现值为左锚点线性插值降尺度到 10 分钟网格，与形状缩放的差异约 1\%。

附件三的 0:00 预报精度反而低于自建持续法（近 7 日同时段均值）：全天光伏电量平均绝对误差分别为 3\,007 与 2\,227 kWh，反方差融合后降至 1\,905 kWh；误差随提前期从 41.4 kW（$h{=}1$）单调劣化至 924.1 kW（$h{=}24$），与提前期无关的持续法在 $h\approx9\sim10$ 交叉。官方预报的价值集中在日内更新而非 0:00。负载侧仍是不确定性大头：日电量平均误差 3\,187 kWh，大于融合后的光伏误差。

\subsection{附件四与前三问的关系}

附件一电价逐时段等于附件四全年 365 日的算术平均（最大偏差 $5.15\times10^{-5}$ 元/kWh，属四舍五入量级）。因此题一至题三相当于用电价的无条件期望做规划，题四改用条件期望与残差情景。附件四全年峰谷比 236.0（附件一为 3.76）；乘性分解 $p_t^{(j)}\approx\text{水平}_j\times\text{形态}_t\times\varepsilon_{jt}$ 解释对数价格的 92.32\% 方差，日内形态高度稳定（逐日与典型形态的相关系数中位数 0.976）。价格呈强周内效应：周五、周六日均价 0.629 元/kWh，其余五天 0.819--0.824 元/kWh，与负载的周内结构完全同相；价格残差与负载残差日级相关系数 0.8869，即“预报低估负载的日子同时也低估了电价”。

% ===================== 五、问题一 =====================
\section{问题一：完美信息下的典型日线性规划}\label{sec:q1}

\subsection{模型建立}

母线是微网内部所有电源与负荷共接的电气节点，自身不储存电能，任一时刻流入必须等于流出。将弃电量 $w_t$ 作为松弛变量，供需平衡写成等式；储能按 SOC 递推演化；昼夜循环约束 $S_{144}=S_0$ 保证“今天占电池的便宜明天还回去”。决策量为计划购电 $g_t$、充电 $c_t$、放电 $d_t$、弃电 $w_t$ 与初始水位 $S_0$（决策化：只要求循环成立且落在运行区间内，最优值由模型给出）：

\begin{equation}
\begin{aligned}
\min\quad & Z=\sum_{t\in\mathcal{T}} p_t\, g_t \\
\text{s.t.}\quad
& g_t+r_t+d_t=\ell_t+c_t+w_t && \forall t \quad\text{（供需平衡）}\\
& S_t=S_{t-1}+\eta\, c_t-d_t/\eta && \forall t \quad\text{（SOC 递推）}\\
& S_{144}=S_0,\quad \underline S\le S_t\le \overline S && \\
& 0\le c_t\le C^{\max},\ 0\le d_t\le D^{\max},\ g_t\ge0,\ w_t\ge0 &&
\end{aligned}
\label{eq:q1}
\end{equation}

共 721 个变量、289 条等式约束的纯线性规划，用 HiGHS 求解器\cite{huangfu2018}求解（亚秒级）。

\subsection{两条结构性质}

\textbf{性质 1（充放天然互斥，模型保持纯 LP）。}若某最优解在时段 $t$ 同时 $c_t=a>0$、$d_t=b>0$，构造新解 $c'_t=a-x$、$d'_t=b-\eta^2x$（$0<x\le\min\{a,b/\eta^2\}$）：SOC 递推项的增量 $\eta(-x)-(-\eta^2 x)/\eta=0$，储能轨迹严格不动，全部边界约束自动保持；平衡式左端减少 $\eta^2x$、右端减少 $x$，富余 $x(1-\eta^2)>0$：若该时段在购电，下调 $g_t$ 使目标严格下降，与最优性矛盾；否则将富余计入 $w_t$，得到等价最优解。逐时段施行至互斥成立。等量缩减构造（$c,d$ 同减）不能保持 SOC 轨迹不变，会沿时间传播破坏循环约束，故必须用 $\eta^2$ 权重的非等量缩减。数值上另建含 0--1 变量与大 M 约束的混合整数规划复解同一问题，最优值与 LP 之差为 0。

\textbf{性质 2（套利门槛 $p_j/p_i\ge1/\eta^2\approx1.2346$）。}设最优解在时段 $i$ 购电充电、时段 $j$ 放电。由互补松弛：充电时段 $\lambda_i=\eta\nu$ 且 $\lambda_i=p_i$，放电时段 $\nu/\eta=\lambda_j\le p_j$，联立得 $p_j/p_i\ge1/\eta^2$。初等验证：为在 $j$ 时段母线多得 1 kWh，需自池内取出 $1/\eta$，故需在 $i$ 时段购 $1/\eta^2$，成本 $p_i/\eta^2$，收益 $p_j$。本题峰谷比 $1.3952/0.3713=3.76$ 远超门槛，套利空间充裕，储能满幅运行，容量约束成为紧约束。该门槛是后三问解释解结构的公共标尺\cite{sioshansi}。

\subsection{求解结果}

最优购电费 \textbf{35\,118.5986 元}，全天购电 59\,482.70 kWh，循环水位 $S_0^\star=S_{144}=8\,550$ kWh。最优调度呈典型的“双谷充、双峰放”形态（图 \ref{fig:q1sched}）：谷段大量购电并同时充电（0--4 时充 1\,500 kWh，20--24 时充 7\,350 kWh），峰段储能放电覆盖净负荷高峰；SOC 轨迹在日内触及上、下界，容量与功率约束成为紧约束，印证性质 2 的判断。

\begin{figure}[H]
\centering
\begin{minipage}[t]{0.485\textwidth}
  \centering \includegraphics[width=0.80\textwidth]{fig04_题1典型日调度.png}
  \caption{问题一最优调度：功率流、SOC 轨迹与储能价值空间（$p_t-\lambda_t$）}\label{fig:q1sched}
\end{minipage}\hfill
\begin{minipage}[t]{0.485\textwidth}
  \centering \includegraphics[width=0.80\textwidth]{fig05_典型日电量对账.png}
  \caption{典型日电量对账：购电与光伏收入对负载与损耗支出，电池侧口径恒等}\label{fig:q1bal}
\end{minipage}
\end{figure}


题目要求的表 1（指定时段计划购电量）与表 2（储能充放电量，电池侧口径）如下：

\begin{table}[H]
\centering\small
\caption{问题一表 1：指定时段计划购电量与全天汇总}
\begin{tabular}{lrrrrrrrr}
\toprule
 & \multicolumn{6}{c}{指定时段购电量（kWh）} & 全天购电量 & 全天购电费 \\
\cmidrule(lr){2-7}
 & 10:00 & 12:00 & 14:00 & 16:00 & 18:00 & 20:00 & （kWh） & （元） \\
\midrule
典型日 & 0.00 & 480.41 & 0.00 & 445.43 & 531.89 & 0.00 & 59\,482.70 & 35\,118.60 \\
\bottomrule
\end{tabular}
\label{tab:q1t1}
\end{table}

\begin{table}[H]
\centering\small
\caption{问题一表 2：储能充放电量（电池侧口径，kWh）}
\begin{tabular}{lrrrrrr}
\toprule
时段 & 0--4 时 & 4--8 时 & 8--12 时 & 12--16 时 & 16--20 时 & 20--24 时 \\
\midrule
充电量 $\eta c$ & 1\,500.00 & 750.00 & 4\,309.17 & 4\,757.43 & 0.00 & 7\,350.00 \\
放电量 $d/\eta$ & 0.00 & 7\,073.16 & 1\,892.22 & 101.22 & 6\,422.37 & 3\,177.63 \\
\bottomrule
\end{tabular}\\[2pt]
{\small 0:00 储电量 $=$ 24:00 储电量 $=$ 8\,550.00 kWh。}
\label{tab:q1t2}
\end{table}

对比基准回答两个层次的问题（图 \ref{fig:q1cost}）：无电池直接购电需 48\,052.05 元，储能带来 \textbf{26.92\%} 节省（要不要电池）；“低于 35\% 价格分位充、高于 65\% 分位放”的阈值规则策略需 36\,712.76 元，优化再省 \textbf{4.34\%}（要不要优化）。能量对账见图 \ref{fig:q1bal}：全天购电 + 光伏 = 负载 + 往返损耗 3\,940.73 kWh，电池侧充放差恰为零，逐项残差为 0。

\begin{figure}[H]
\centering
\begin{minipage}[t]{0.485\textwidth}
  \centering \includegraphics[width=0.62\textwidth]{fig06_题1费用对照.png}
  \caption{问题一费用对照：优化解、阈值规则与无电池直买}\label{fig:q1cost}
\end{minipage}\hfill
\begin{minipage}[t]{0.485\textwidth}
  \centering \includegraphics[width=0.62\textwidth]{fig07_S0敏感性.png}
  \caption{初始水位 $S_0$ 扫描：费用带宽 0.24\%，循环水位近乎自由方向}\label{fig:q1s0}
\end{minipage}
\end{figure}


\subsection{对偶分析与模型检验}

对偶变量的读数与理论一致（图 \ref{fig:q1dual}）：购电时段母线边际价值 $\lambda_t$ 与电价 $p_t$ 严格重合，放电时段 $\lambda_t$ 明显低于 $p_t$（电池成为更便宜的边际电源）；储能边际价值 $\nu_t$ 全天平滑变化，循环约束对偶 $\nu^\star=0.4720$ 元/kWh 给出“池内留 1 kWh 过夜”的边际价值，将成为题二、题三的终端价值锚。自检：循环误差 $|S_{144}-S_0|=0$，逐时段弃电为 0（全额消纳），充放乘积为 0。灵敏度方面（图 \ref{fig:q1s0}），把 $S_0$ 逐点固定在 1\,200 至 10\,800 kWh 重解，费用曲线在 8\,550 附近呈平坦盆地、总带宽仅 0.24\%；特别地，固定为附录一标定的 6\,000 kWh 得 $Z^{6000}=35\,126.95$ 元，仅比自由解贵 8.35 元（$+0.024\%$）。该扫描同时给出两个后文要用的结论：其一，$S_0$ 决策化不损害稳健性，把循环水位交给模型是安全的；其二，循环水位在费用上近乎退化，因此“日内 SOC 均值”不具判定力，后文比较各问的储能保有量时一律改用\textbf{隔夜水位}。外生量的完整清单与终端价值的标定实验见附录 B。

\begin{figure}[H]
\centering
\includegraphics[width=0.80\textwidth]{q1/q1_duals.png}
\caption{问题一对偶变量：母线边际价值 $\lambda_t$ 与储能边际价值 $\nu_t$ 对电价 $p_t$}
\label{fig:q1dual}
\end{figure}

% ===================== 六、问题二 =====================
\section{问题二：滚动随机优化与在线随机 MPC}\label{sec:q2}

\subsection{两阶段决策结构}

在假设 A5（诚实回测）下，第 $j$ 日 0:00 可见信息集 $\mathcal{F}_0^{(j)}$ 只含历史。决策分两阶段：计划购电量 $g^{(j)}$ 于 0:00 锁定、全额计费、不可增减；充放、紧急购电、弃电为实时补救量。全年以真实路径 SOC 链衔接（$S_0^{(j+1)}=S_{144}^{(j)}$），一月暖机不计费。整体架构见图 \ref{fig:q2arch}。

\begin{figure}[H]
\centering
\includegraphics[width=0.80\textwidth]{fig09_题2四层架构与随机MPC.png}
\caption{问题二四层架构与在线随机 MPC：当步决策确定、剩余时域按情景分支}
\label{fig:q2arch}
\end{figure}

\subsection{预报层：情景生成}

负载点预报取最近 4 个同星期几日的同点均值（居民负载的周节律），光伏取近 7 日同点均值（天气持续性与慢季节漂移，无周节律）。误差情景采用“残差整条贴加”：对每个历史日，用其当时的真预报（伪样本外，无信息泄漏）计算负载加性残差与光伏相对残差，取最近 12 条残差日曲线、整条贴加到当日点预报上，等概率构成 $N_s=12$ 个情景。整条贴加自动保留日内相关与负荷—光伏联动，无需显式估计协方差矩阵。两个窗口 $W_L=4$、$W_R=7$ 的去超参验证见附录 \ref{sec:appBwin}；负载/光伏/残差窗的 $3\times3\times2$ 选臂网格见第 \ref{sec:aux} 节。交付口径下负载预报平均绝对误差 31.48 kWh/时段、光伏 25.57 kWh/时段，误差分布与季节结构见图 \ref{fig:q2fcst}。

\begin{figure}[H]
\centering
\includegraphics[width=0.80\textwidth]{fig10_题2预报质量-L1P1.png}
\caption{问题二交付组合的预报质量：负载与光伏的误差分布及分季节技能}
\label{fig:q2fcst}
\end{figure}

\subsection{规划层：日前两阶段场景线性规划}

计划 $g$ 不带情景上标，即一个计划服务全部情景（非预期性约束的紧凑写法），这是两阶段随机规划区别于逐情景优化的本质\cite{birge}：

\begin{equation}
\begin{aligned}
\min_{g\ge0}\quad
& \tilde Z^{(j)}=\sum_{t}p_t g_t^{(j)}
+\sum_{s\in\mathcal S^{(j)}}\pi_s\Big[\kappa\sum_{t}p_t u_t^{(j,s)}-\nuend S_{144}^{(j,s)}\Big]\\
\text{s.t.}\quad
& g_t+d_t^{(s)}+u_t^{(s)}=n_t^{(s)}+c_t^{(s)}+w_t^{(s)},\quad
S_t^{(s)}=S_{t-1}^{(s)}+\eta c_t^{(s)}-d_t^{(s)}/\eta,\quad\forall s
\end{aligned}
\label{eq:q2plan}
\end{equation}

末端估值 $-\nuend S_{144}$ 消除滚动时域的末端效应（否则最优解在日末抽干电池、把成本转嫁次日）。$\nuend$ 的经济学区间为 $[p^{\mathrm v}/\eta,\ \eta p^{\mathrm p}]=[0.4126,1.2557]$ 元/kWh（重置成本与变现收益），区间非空当且仅当峰谷比不低于套利门槛 $1/\eta^2$，与问题一的理论对接；锚定值取问题一循环约束对偶 0.472 元/kWh。该量是规划层唯一需要标定的人工参数，三重验证见附录 \ref{sec:appBnu}。

\subsection{紧急购电的口径上界}\label{sec:q2cap}

题面语义中紧急购电是断供兜底。若只把它写成“无上界、按 $\kappa p_t$ 计费”的自由变量，它将额外获得日内采购通道的身份：计划锁定后日内增加入池电量只剩此一条路，而平价段 5 倍价买入折损后约 4.44 元/kWh，比晚峰 5 倍价补缺口（6.98 元/kWh）便宜，模型就会用它囤货。本文早期版本实测 41.7\% 的紧急电量属囤电型（发生在电池不空、当步仍在充电的时段），紧急购电量的 9.51\% 甚至落在光伏盈余时段。这是可行域给多了，不是求解器出错。

修正为口径上界 $u_t\le(n_t-g_t)^+$：右端由数据与已锁定的 $g$ 完全决定，零新参数；可证其等价于禁止“紧急购电同时充电”（由平衡式与充放互斥：$u_t>0$ 蕴含 $n_t>g_t$ 且 $c_t+w_t\le d_t$，再由互斥性迫使 $c_t=0$）；恒可行性构造 $c=d=0$、$u=(n-g)^+$、$w=(g-n)^+$ 恰取在上界上，全年回测不会中断。右端含第一阶段决策 $g$，作为 $g$ 的约束非凸，采用序列线性化迭代到不动点：每轮用上一轮的 $g$ 固定上界、重解纯 LP，直至 $\|\Delta g\|_1/\sum g\le10^{-6}$。实测 331/334 天机器精度收敛（中位数 2 轮），其余 3 天为等价最优之间的缓慢棘轮，按经济收敛判过（末轮目标变化 $\le9.3\times10^{-6}$）。

\subsection{执行层：在线随机 MPC}

每个 10 分钟时段解一个“当步 here-and-now + 剩余时域按情景分支”的两阶段子问题：当步决策与 SOC 在全部情景间共享、当步用真实净负荷与真实 SOC，只执行本时段。保留情景分布的必要性可由影子价格看出：设剩余时域某峰时段以概率 $q$ 出现缺口，池内 1 kWh 的影子价格为 $q\cdot\kappa p^{\mathrm p}\eta$；当步小缺口用放电优于紧急购电当且仅当 $q<p_t/(p^{\mathrm p}\eta^2)$。确定性等价（把情景集压缩成单条分位参考路径）相当于令 $q\equiv1$，谷价时段门槛为 0.385 恒不满足，“电池有电也不放、改在谷时以 5 倍价买”成为其最优解。随机 MPC 把 $q$ 还原为真实频率，判据恢复正常；同时人工分位参数被消去。

\subsection{求解结果}

\begin{table}[H]
\centering\footnotesize
\caption{问题二全年结果与信息阶梯（334 个交付日，元）}
\begin{tabular}{lrr}
\toprule
口径 & 费用 & 说明 \\
\midrule
$\ZPI$ 完美信息 & 12\,223\,092.57 & 0:00 已知当日真值，理论下界 \\
$\Zep$ 事后最优补救 & 13\,463\,630.76 & 给定 $g$ 的执行极限 \\
$\Zon$ 在线随机 MPC（交付） & \textbf{13\,745\,816.13} & \\
$\ZCE$ 确定性等价 MPC & 14\,254\,554.72 & 同 $g$、同 SOC，只差未来表示 \\
\midrule
无电池（报童最优 $\alpha{=}0.85$） & 17\,096\,785.82 & 相对交付 $+24.4\%$ \\
谷充峰放规则策略 & 16\,010\,662.24 & 相对交付 $+16.5\%$ \\
问题一计划逐日复读 & 24\,816\,066.97 & 相对交付 $+80.5\%$ \\
\bottomrule
\end{tabular}
\label{tab:q2ladder}
\end{table}

全年 334 个交付日合计费用 \textbf{13\,745\,816.13 元}（表 \ref{tab:q2ladder}），由两笔构成：计划购电费 12\,818\,726 元（93.3\%）与紧急购电费 927\,090 元（6.74\%）——5 倍罚则只作用于不到百分之七的支出，计划的总体覆盖水平恰当。分日看，夏至日（6.21）光伏覆盖近半负载、全天仅购 35\,318 kWh；冬至日（12.21）净负荷最高、购 95\,949 kWh（表 \ref{tab:q2days}）；紧急购电只在春分、秋分两个指定日发生，秋分日午后至晚峰共 7 段 2\,259 kWh（表 \ref{tab:q2emerg}），对应连日阴雨后的预报偏差期。

\begin{table}[H]
\centering\footnotesize
\caption{问题二表 1：四个指定日的计划购电量（kWh）与购电费（元，含紧急费）}
\begin{tabular}{lrrrrrrrr}
\toprule
 & \multicolumn{6}{c}{指定时段计划购电量} & 全天购电量 & 全天购电费 \\
\cmidrule(lr){2-7}
日期 & 10:00 & 12:00 & 14:00 & 16:00 & 18:00 & 20:00 & （kWh） & （元） \\
\midrule
03-20（春分） & 0.00 & 682.93 & 0.00 & 472.59 & 681.32 & 0.00 & 67\,189.46 & 42\,111.37 \\
06-21（夏至） & 0.00 & 0.00 & 0.00 & 144.28 & 189.05 & 0.00 & 35\,318.42 & 20\,651.78 \\
09-23（秋分） & 0.00 & 435.19 & 0.00 & 491.94 & 721.28 & 17.80 & 64\,084.53 & 52\,704.62 \\
12-21（冬至） & 0.00 & 958.12 & 0.00 & 840.21 & 711.94 & 0.00 & 95\,949.37 & 61\,723.87 \\
\bottomrule
\end{tabular}
\label{tab:q2days}
\end{table}

\begin{table}[H]
\centering\small
\caption{问题二表 3：四个指定日的紧急购电情况}
\begin{tabular}{lrp{5.6cm}rr}
\toprule
日期 & 段数 & \multicolumn{1}{c}{购电时间段} & 紧急购电量（kWh） & 紧急购电费（元） \\
\midrule
03-20 & 3 & 17:00--17:20、20:10--20:40、21:30--21:50 & 226.22 & 1\,322.78 \\
06-21 & 0 & 未发生 & 0.00 & 0.00 \\
09-23 & 7 & 14:50--15:00、15:10--15:20、15:50--16:50、18:20--18:50、19:10--19:20、19:40--20:00、20:10--21:50 & 2\,259.14 & 12\,764.51 \\
12-21 & 0 & 未发生 & 0.00 & 0.00 \\
\bottomrule
\end{tabular}
\label{tab:q2emerg}
\end{table}

在总数之外，费用还能拆出“信息的价值”：完美信息下界 $\ZPI=12\,223\,093$ 元，$\mathrm{EVPI}=152.3$ 万元是“不知道未来”的总代价，其中 81.5\% 丢在 0:00 的计划上——改进优先级应在预报与计划层；同一计划、同日初 SOC 下把执行器换成确定性等价（表 \ref{tab:q2ladder} 末行）则贵 50.9 万元，这是“未来以分布还是单条路径进入模型”的净价差（图 \ref{fig:q2ladderattr}）。表下半三行是策略类副模型（无电池、规则、题一复读），完整归因纪律与四问信息阶梯见第 \ref{sec:aux} 节。逐日费用与紧急费的时间分布见图 \ref{fig:q2daily}（紧急费集中在负载高、光伏弱的冬季与早春），四个指定日的剖面见图 \ref{fig:q2solstice}。

\begin{figure}[H]
\centering
\begin{minipage}[t]{0.485\textwidth}
  \centering \includegraphics[width=0.80\textwidth]{fig11_题2逐日费用与紧急费.png}
  \caption{问题二逐日总费用与紧急购电费：季节模式与异常日}\label{fig:q2daily}
\end{minipage}\hfill
\begin{minipage}[t]{0.485\textwidth}
  \centering \includegraphics[width=0.80\textwidth]{fig12_题2二分二至剖面.png}
  \caption{问题二四个指定日（春分、夏至、秋分、冬至）的购电与 SOC 剖面}\label{fig:q2solstice}
\end{minipage}
\end{figure}


紧急购电结构：发生 168/334 天、2\,066 个时段，频率 4.30\%，与报童临界频率 $1-\alpha^\star=4.5\%$（过量入池回收率 $\theta=\eta^2$）一致\cite{qin2011}，该频率是计划层覆盖水平的设计值而非失控；电量按电价分段为谷 5.57\%、平 42.85\%、峰 51.58\%（真缺口集中在峰时），“紧急购电同时充电”的时段数为 0，光伏盈余时段占比 0.00\%。归因五分类（按电量）：放电功率强制 8.71\%、储能耗尽 18.99\%、囤电型 0.00\%（结构保证）、配给型 49.06\%、低水位带 23.23\%。配给型指情景集认为晚峰缺口概率高于门槛时，把有限存量留给更贵窗口、当步平价缺口用紧急购电顶替，属最优跨时段配给而非失误（抽检 24 个此类日，全部在事后把 SOC 降至 2\,500 kWh 以下且均在电价高于 1.2 元的时段放电）。弃电集中于 9--14 时：该时段容量与充电功率双紧（48.5\% 的弃电时段、11.1\% 时段充电顶格），是“存不下”而非“不想存”。

值得说明的是口径上界的取舍。加与不加上界，全年费用仅相差 0.17\%（23\,338 元），但交付表的含义完全不同：无上界时模型会在电池接近满电时仍以 5 倍价买入囤存，这在数学上最优（\ref{sec:q2cap} 节的套利机制使然），但“一边宣称断供紧急、一边在电池充盈时高价囤电”对小区居民难以解释，舆情风险与赔付暴露都更高。本文放弃这 0.17\% 的费用优势，换取交付表中每一度紧急购电都对应一次真实断供——这是用可解释性定价的合规成本。

\begin{figure}[H]
\centering
\includegraphics[width=0.80\textwidth]{fig13_题2费用阶梯与紧急归因.png}
\caption{问题二费用阶梯（左）与紧急购电归因的加/不加口径上界对照（右）}
\label{fig:q2ladderattr}
\end{figure}

\subsection{模型检验}

\textbf{执行层与口径。}固定 MPC 与随机 MPC 的配对对照、以及紧急购电上界的合规成本，统一放在第 \ref{sec:aux} 节展开；此处只保留本问特有的三代数字：把未来压成单条分位路径时年费用 14\,075\,337 元（谷价紧急购电 51\%、缺口频率 11.19\%、86.5\% 紧急电量发生在 SOC 高于 3\,000 kWh）；换成随机 MPC 后 13\,722\,478 元；再加口径上界得交付 13\,745\,816 元，合规成本 $+0.17\%$。

\textbf{规划层已在费用最优点的证据。}这组实验回答“紧急购电多，是不是计划买少了”：把 12 个情景的净负荷人为整体上抬 50 kWh/时段（相当于强迫计划每天多买约 7\,200 kWh），其余不动、全年重跑。结果紧急电量确实下降 71\%（指标变好），但总费用上升 4.79\%、弃电增加 65\%（中午容量与充电功率双紧，多买的电存不下只能弃掉）。它证伪了“紧急购电多 $=$ 计划不足”的直觉：计划已在费用最优点，继续加厚是以钱买心安。末端估值的局部重扫与全年地标扫描见附录 \ref{sec:appBnu}：锚 0.472 两侧均变贵。

\textbf{自检。}对账恒等式（充放电池侧差 $=$ SOC 日内净变化）与能量守恒残差 $\le10^{-11}$；序关系 $\ZPI\le\Zep\le\Zon$ 全年 0 违例；充放乘积与紧急—弃电乘积逐时段为 0。

% ===================== 七、问题三 =====================
\section{问题三：分块承诺与预报融合}\label{sec:q3}

\subsection{结算规则的化简与分块承诺}

记计划 $g_t$、调整 $a_t$，偏差 $\delta^\pm=(\pm(a_t-g_t))^+$。三段计价（重合部分原价、下调按 $\beta^-=0.5$、上调按 $\beta^+=1.5$）化简为
\begin{equation}
Z_t^{\mathrm{buy}}=p_t\min(g_t,a_t)+\beta^-p_t\delta_t^-+\beta^+p_t\delta_t^+
=p_ta_t+\tfrac12 p_t|a_t-g_t|,
\label{eq:q3settle}
\end{equation}
化简只用到 $\beta^+-1=\beta^-=0.5$：看似不对称的两向罚则实为一条对称的 $L_1$ 惩罚。$g_t$ 自 0:00 起只是计费基准，供需平衡中出现的是 $a_t$；令 $a\equiv g$ 则逐字退化回问题二，该特例关系是本题的复现检验。

\textbf{块分解定理。}每个时段 $t$ 的调整只在末次修订机会 $\tau_{m(t)}$ 执行（0:00 报全天 $g$；6 时定 6--12 时块、12 时定 12--18 时块、18 时定 18--24 时块；0--6 时块永不可改），弱占优于任何更早或反复修订：链式结算下由三角不等式
$\tfrac12 p_t(|a^{(1)}-g|+|a^{(2)}-a^{(1)}|)\ge\tfrac12p_t|a^{(2)}-g|$，
且“费用挂时段电价、信息单调、远期承诺不约束中间执行”三支柱与结算读法无关。该定理把 $4\times144$ 维修订决策约简为四个 36 维块决策，是全年 334 天可解的前提。分块结构与四档价签（违约 0.5 / 计划 1.0 / 加订 1.5 / 紧急 5）见图 \ref{fig:q3commit}。

\begin{figure}[H]
\centering
\includegraphics[width=0.80\textwidth]{fig14_题3分块承诺与四档价签.png}
\caption{题三分块承诺时间结构与四档价签阶梯}
\label{fig:q3commit}
\end{figure}

\subsection{预报层：反方差融合与日内缩放}

光伏点预报取官方预报与持续法的凸组合，权重按两估计量的误差方差反比确定\cite{bates1969}：
\begin{equation}
\varphi(h)=\frac{\sigma_{\mathrm{pe}}^2(h)}{\sigma_{\mathrm{pe}}^2(h)+\sigma_{\mathrm{fc}}^2(h)},\qquad
\hat R_t=\varphi(h_t)\,\tilde R_t+\bigl(1-\varphi(h_t)\bigr)\,\bar R_t,
\end{equation}
方差只用当日之前的数据滚动估计。短提前期官方预报占优（$h{=}1$ 权重 0.98），长提前期持续法占优（$h{=}24$ 权重 0.11），理论权重与实测最优权重吻合（图 \ref{fig:q3fusion}），融合是推导结果而非调参。负载预报继承问题二，并在各修订时刻按当日已实现段的比例缩放更新（截断 $[0.7,1.3]$，午后与晚间块的平均误差分别下降 39\% 与 32\%）。误差情景按修订阶段与提前期区间分池抽取。

\begin{figure}[H]
\centering
\begin{minipage}[t]{0.485\textwidth}
  \centering \includegraphics[width=0.80\textwidth]{fig15_题3白天预报融合.png}
  \caption{题三白天目标预报：提前期 RMSE（官方 / 持续法 / 融合）与反方差权重 $\varphi(h)$}\label{fig:q3fusion}
\end{minipage}\hfill
\begin{minipage}[t]{0.485\textwidth}
  \centering \includegraphics[width=0.80\textwidth]{fig16_题3修订演示-最大偏差日.png}
  \caption{最大偏差日的修订演示：计划 $g$、调整 $a$ 与块锁定时刻}\label{fig:q3demo}
\end{minipage}
\end{figure}


数据核验给出的结构判断贯穿建模：光伏 99.5\% 集中在 6--18 时，18--24 时几乎无光伏却是全天最贵的块（均价 0.926 元/kWh），故 18:00 修订的信息价值主要来自已观测的负载而非光伏预报；0--6 时恒为最便宜的块。

\subsection{规划与执行：嵌套双环}

0:00 沿用问题二的两阶段场景规划给出全天计划 $g$（$g$ 就是打算执行的日前购电；不做“假定后面能改”的松弛（该松弛缺失非预期性，实测年费用贵 15.1 万元））。6/12/18 时解分块两阶段子问题：本块调整量 $a_{\mathcal{B}_m}$ 是 here-and-now（不带情景上标、收偏差费），更晚块在情景内试探；约束沿用共用物理可行域，$g$ 换为 $a$，紧急购电上界同步为 $u_t\le(n_t-a_t)^+$。内环为问题二同款在线随机 MPC（无权改 $a$），外环锁定值与真实 SOC 进入下一阶段。全程末端估值同 0.472。

\subsection{求解结果}

\begin{table}[H]
\centering\small
\caption{问题三全年结果（334 个交付日）}
\begin{tabular}{lrr}
\toprule
指标 & 数值 & 对比 \\
\midrule
$\Zon$（交付） & \textbf{13\,023\,147.06 元} & 较问题二省 722\,669 元（$5.25\%$） \\
\quad 基础购电费 $\sum pa$ & 12\,608\,890.51 元 & \\
\quad 偏差费 $\tfrac12\sum p|a-g|$ & 295\,233.33 元 & 占 $2.27\%$ \\
\quad 紧急购电费 & 119\,023.22 元 & 占 $0.91\%$（问题二为 $6.74\%$） \\
$\ZCE$ / 执行层增益 & 13\,106\,182.26 / 83\,035 元 & \\
$\Zep$ / $\ZPI$ & 12\,983\,832.24 / 12\,243\,666.87 & \\
实时信息损失 & 39\,315 元 & 执行层已接近事后最优 \\
弃电率 / SOC 均值 / 隔夜水位 & $6.53\%$ / 6\,988 / 8\,025 kWh & \\
调整上调 / 下调电量 & 277\,085 / 522\,388 kWh & 行使率 $22.9\%$ \\
\bottomrule
\end{tabular}
\label{tab:q3res}
\end{table}

按修订机会把全天分为四块：0--6 时块在 0:00 报出计划后即成最终量（其后没有修订机会，称不可修订块）；6--12、12--18、18--24 时三个块分别可在 6/12/18 时修订，称可修订块。问题三的结构变化全部兑现（表 \ref{tab:q3res}、图 \ref{fig:q3demo} 与图 \ref{fig:q3q2}）：1.5 倍加订通道顶替了 5 倍兜底，紧急费占比从 6.74\% 压至 0.91\%；分块对比问题二，三个可修订块的紧急电量合计降至其 0.126 倍（调整通道吸收了缺口），不可修订的夜间块为 0.003 倍（其缺口风险本与问题二同级，压降来自更准的日前计划）；计划覆盖水平在 6:00 处出现台阶（夜间块覆盖率 1.038，其后三块均值 0.930），即夜间无修订机会故多备、白天可修订故少备。偏差费 2.27\% 是购买修订灵活性的定价。下调电量（52.2 万 kWh）多于上调（27.7 万 kWh），全部落在三个可修订块（夜间不可修订块为 0），其中午后两块占下调的 90.3\%；机理是日内预报把光伏与负载向下修正、原计划偏多，与“电池有余量时应宁存不退”的理论预期方向相反，本文如实报告该差异。最大偏差日的修订行为演示见图 \ref{fig:q3demo}，题二与题三的逐日费用对照见图 \ref{fig:q3q2}：5.25\% 的净节省主要来自紧急费的压降，部分被偏差费抵消。

四个指定日的计划与调整对照见表 \ref{tab:q3days}：春分、夏至的调整量低于计划（日内预报下修），秋分则上修，冬至基本不动——调整方向与当日预报偏差一致。

\begin{table}[H]
\centering\footnotesize\setlength{\tabcolsep}{3pt}
\caption{问题三：四个指定日的计划与调整购电量（kWh）及全天汇总}
\begin{tabular}{lrrrrrrrrrrrrrr}
\toprule
 & \multicolumn{6}{c}{计划购电量 $g$} & \multicolumn{6}{c}{调整购电量 $a$} & \multicolumn{2}{c}{全天量 $g/a$} \\
\cmidrule(lr){2-7}\cmidrule(lr){8-13}\cmidrule(lr){14-15}
日期 & 10 & 12 & 14 & 16 & 18 & 20 & 10 & 12 & 14 & 16 & 18 & 20 & （kWh） & 购电费（元） \\
\midrule
03-20 & 0 & 543 & 0 & 553 & 681 & 0 & 0 & 543 & 0 & 553 & 686 & 0 & 65\,821/65\,763 & 42\,019 \\
06-21 & 0 & 0 & 0 & 132 & 37 & 0 & 0 & 0 & 0 & 95 & 0 & 0 & 35\,638/35\,066 & 20\,262 \\
09-23 & 0 & 336 & 0 & 492 & 721 & 18 & 0 & 289 & 0 & 492 & 745 & 18 & 64\,712/65\,963 & 45\,826 \\
12-21 & 0 & 995 & 0 & 850 & 709 & 0 & 0 & 934 & 0 & 860 & 685 & 0 & 96\,024/94\,710 & 61\,823 \\
\bottomrule
\end{tabular}
\label{tab:q3days}
\end{table}

\textbf{附加问：是否需要引入其他时刻的预报。}在合格基准（非预期性计划 + 冻结不调整）上逐级开放修订权：仅 0:00 报计划 13\,633\,398.76 元，开放 6:00 后降至 13\,345\,217.23 元（省 28.8 万元），再开放 12:00 降至 13\,176\,494.13 元（省 16.9 万元），全开后为交付值 13\,023\,147.06 元（再省 15.3 万元）。四级修订合计 \textbf{61.03 万元（4.48\%）}，其中 6:00 档贡献最大（图 \ref{fig:q3ladder}）。结论：需要引入 6/12/18 时的预报与修订，且越早的修订价值越大；18:00 档的信息来源主要是当日已观测负载（附件三只预报光伏）。该阶梯在波动价下的同构对照见第 \ref{sec:aux} 节：修订权合计 $4.55\%$，6:00 档仍近半，说明附加问的答案不依赖价格环境。

\begin{figure}[H]
\centering
\begin{tikzpicture}
\begin{axis}[
  width=0.68\textwidth, height=4.6cm,
  ybar, bar width=26pt,
  symbolic x coords={仅0:00, +6:00, +12:00, 全开(交付)},
  xtick=data,
  ymin=1290, ymax=1380,
  ylabel={全年费用（万元）},
  nodes near coords={\pgfmathprintnumber\pgfplotspointmeta},
  every node near coord/.append style={font=\scriptsize},
  nodes near coords align={vertical},
]
\addplot coordinates {
  (仅0:00,1363.3399) (+6:00,1334.5217) (+12:00,1317.6494) (全开(交付),1302.3147)
};
\end{axis}
\end{tikzpicture}
\caption{修订阶梯：逐级开放 6/12/18 时修订的全年费用（万元）。逐级增量 $-28.8/-16.9/-15.3$ 万元}
\label{fig:q3ladder}
\end{figure}

\begin{figure}[H]
\centering
\includegraphics[width=0.80\textwidth]{fig17_题2与题3费用对照.png}
\caption{题二与题三费用对照：净节省 $5.25\%$ 的构成（紧急费压降 $-$ 偏差费新增）}
\label{fig:q3q2}
\end{figure}

\section{问题四：波动电价下的三元不确定性问题}\label{sec:q4}

\subsection{信息认定与可行性依据}

附件四与附件二同为按日解锁的实现值，故当日电价在 0:00 不可知，需自建预报（假设 A5 的平移）。该读法可行的数据依据：乘性分解解释对数价格 92.32\% 的方差，日内形态高度稳定；0--6 时为当日最便宜块的判定，真值 365/365 成立、预报同样 365/365，调度骨架“夜里囤、傍晚放”不依赖当日价格信息。理论依据：设当日价格为日水平与形态的乘积，可证日内最优解集与日水平无关，仅通过跨日价格比影响日末水位；而形态误差（5.58\%）与水平误差（4.03\%）中真正伤害调度的是前者。

\subsection{价格预报与三元联合情景}

价格点预报取近 4 个同星期几同时段均值（MAPE 7.51\%），技能全部来自周内效应（不含星期结构的 7 日均值 MAPE 13.65\%）。情景生成扩为三元：负载、光伏、价格的残差取自同一历史日、整条贴加。理论依据是一行等式：对第一阶段决策 $x_t$（不依赖情景），$\mathbb{E}_s[p_t^{(s)}x_t]=\hat p_t x_t$ 精确折叠；而对情景相关的补救量 $y_t^{(s)}$，
\begin{equation}
\mathbb{E}_s\bigl[p_t^{(s)}y_t^{(s)}\bigr]=\hat p_t\,\mathbb{E}_s[y_t^{(s)}]+\mathrm{Cov}_s\bigl(p_t^{(s)},y_t^{(s)}\bigr).
\label{eq:cov}
\end{equation}
价格不确定性的全部影响集中在协方差通道。价格与负载残差日级相关 0.8869，缺口大的情景同时是价格高的情景，$\mathrm{Cov}(p,u)>0$：若将价格残差独立抽样，会系统性低估“高缺口 $\times$ 高价”的最坏组合，从而低估紧急购电费用。三元联合贴加使该相关自动保留。

\subsection{终端价值逐日化与门槛重判}

常数终端价值在便宜日囤积过度、昂贵日抽干过度，改为逐日标定：以次日预报价曲线配典型日负载光伏解问题一型循环线性规划，取循环约束对偶的绝对值，全链路不含拟合常数（首日在附件一价格下复现 0.472，为标定自检）。该设定给出可检验预言：周四的次日是低价的周五，日末应降水位；周日的次日是常规周一，日末应抬水位。门槛体系需逐日重判：5 倍罚则下的“任意时刻预防性充电”门槛 $\kappa\eta^2=4.05$，在常数价下全年成立，波动价下仅 113/365 天成立，保险从无条件退化为“提前在当日的便宜块里囤好”；报童临界分位则由物理测度退化为价格加权测度下的分位（缺口与价格同向时最优覆盖水平应更高）。

\subsection{弃电拆项与报童保险}

波动价下平衡松弛 $w_t$ 需要拆开解读：一部分是弃光 $w_{\mathrm{pv}}=\min\{w,(r-\ell)^+\}$（光伏盈余存不下），另一部分是超买溢出 $w_{\mathrm{over}}$（计划锁定后供给超过需求的物理必需）。后者按实现价入账为“报童保险费”，是为压低 5 倍缺口风险而多买的量，不是可省的浪费：报告期中日计划量超过日净负荷的天数占比（日频超买占比）应落在报童区间 $1-1/\kappa=80\%$ 附近，这是 $\kappa=5$ 分位的可直接观测量。序关系检验同步改为末端估值调整口径 $\tilde Z=Z^{\mathrm{cash}}-\nuend S_{144}$：现金口径在波动价下会把完美信息的合理囤库误判为违例（便宜电买成库存推高现金、库存项补偿后仍最优），调整口径消除该假阳性。

\subsection{求解结果}

按题意分两条交付线：Q4-2（问题二骨架 + 波动价）与 Q4-3（问题三骨架 + 波动价），共用物理内核与执行层，各自独立滚动，全年 334 天。主结果见表 \ref{tab:q4res}：调整线较计划线再省 775\,605 元（5.37\%），与常数价下的节省幅度（5.25\%）一致，说明修订机制的价值在波动价下保持；两条线的紧急购电占比分别为 6.62\% 与 0.87\%，峰段占比 71.07\% 与 61.63\%（真缺口集中在贵时段），归因四条判定全部通过。

\begin{table}[H]
\centering\small
\caption{问题四全年结果（读法 R，334 个交付日，元）}
\begin{tabular}{lrr}
\toprule
指标 & Q4-2（计划线） & Q4-3（调整线） \\
\midrule
$\Zon$ & 14\,447\,039.68 & \textbf{13\,671\,435.06}（省 $5.37\%$） \\
$\ZCE$ / 执行层增益 & 14\,771\,011.74 / 323\,972.06 & 13\,708\,593.11 / 37\,158.06 \\
$\Zep$ & 14\,137\,808.93 & 13\,635\,251.00 \\
$\ZPI$（给定本臂 $S_0$） & 12\,707\,701.21 & 12\,741\,991.88 \\
紧急购电费 / 占比 & 955\,821.96（$6.62\%$） & 119\,323.31（$0.87\%$） \\
紧急峰段占比 / 发生天数 & 71.07\% / 165 天 & 61.63\% / 133 天 \\
弃光率（分母 $=$ 母线盈余） & 44.11\% & 34.22\% \\
超买溢出占计划 / 报童保险费 & 2.08\% / 327\,894.83 元 & 0.90\% / 109\,108.52 元 \\
日频超买占比（理论 $80\%\pm10\%$） & 88.62\%（过） & 86.23\%（过） \\
SOC 均值 / 日末（kWh） & 6\,177.6 / 10\,300.0 & 5\,951.3 / 9\,731.4 \\
\bottomrule
\end{tabular}
\label{tab:q4res}
\end{table}
计划线费用较问题二（13\,745\,816 元）高 701\,224 元（$5.1\%$）。这个差额必须先拆开才能解读，否则会被误读成“波动价下模型变差了”：把问题二的交付轨迹原样搬到附件四实现价下重开账单（不重新优化）得到中间量 $\Zmis=14\,519\,354$ 元，于是 $5.1\%$ 分成\textbf{价格效应} $+773\,538$ 元与\textbf{策略适应} $-72\,314$ 元两格——贵出来的钱全部来自价格环境（附件一是附件四的全年均值，而高负荷日恰是高价日，按无条件期望规划再按实现价结算必然偏贵），而模型对波动价的适应反过来挣回 7.2 万元。调整线同构，两格为 $+721\,636$ 与 $-73\,348$ 元。该分解与价格消融、信息上界一并归入副模型体系（\ref{sec:aux} 节），此处只取结论：\textbf{题四不可能比题二便宜，但在同一价格环境下本文的策略优于把常数价方案照搬}。

消融给出的三笔账见图 \ref{fig:q4ablate}（均为 334 天、计划线）：不建模价格（价格预报换附件一典型价）贵 38\,948 元，终端价值退回常数 0.472 贵 19\,664 元，若 0:00 即知当日全部实现价（读法 D）省 44\,762 元。电价的“可学到部分”与“当日才知道的部分”价值量级相当（各约 $0.3\%$），与数据核验中形态误差 $5.58\%$ 对水平误差 $4.03\%$ 的判断一致；两笔账的严格分离见 \ref{sec:aux} 节。

\begin{figure}[H]
\centering
\begin{minipage}[t]{0.485\textwidth}
  \centering \begin{tikzpicture}
\begin{axis}[
  width=0.78\textwidth, height=4.9cm,
  ybar, bar width=22pt,
  symbolic x coords={不建模价格, 常数终端价值, 主交付(读法R), 知当日价(读法D)},
  xtick=data, x tick label style={font=\scriptsize},
  ymin=1436, ymax=1452,
  ytick={1438,1440,1442,1444,1446,1448,1450,1452},
  ylabel={全年费用（万元）},
  nodes near coords={\pgfmathprintnumber\pgfplotspointmeta},
  every node near coord/.append style={font=\scriptsize},
  nodes near coords align={vertical},
]
\addplot coordinates {
  (不建模价格,1448.5988) (常数终端价值,1446.6703) (主交付(读法R),1444.7040) (知当日价(读法D),1440.2278)
};
\end{axis}
\end{tikzpicture}
  \caption{问题四价格信息与建模价值阶梯（计划线，万元）：不建模价格最贵，知当日实现价最省}\label{fig:q4ablate}
\end{minipage}\hfill
\begin{minipage}[t]{0.485\textwidth}
  \centering \begin{tikzpicture}
\begin{axis}[
  width=0.7\textwidth, height=4.8cm,
  ybar, bar width=26pt,
  symbolic x coords={次日周五, 次日周六, 次日周一, 其他日期},
  xtick=data,
  ymin=0, ymax=8600,
  ylabel={日末水位（kWh）},
  nodes near coords={\pgfmathprintnumber\pgfplotspointmeta},
  every node near coord/.append style={font=\scriptsize},
  nodes near coords align={vertical},
]
\addplot coordinates {
  (次日周五,2594) (次日周六,6480) (次日周一,6309) (其他日期,6974)
};
\end{axis}
\end{tikzpicture}
  \caption{计划线日末水位的星期台阶（性质 9）：次日为周五时主动清仓}\label{fig:q4week}
\end{minipage}
\end{figure}

四个指定日的结果见表 \ref{tab:q4d2} 与表 \ref{tab:q4d3}（费用按实现价结算，为事后量；两线的计划量不同，因调整线知道计划可被修订而报得更克制）。

\begin{table}[H]
\centering\footnotesize\setlength{\tabcolsep}{3pt}
\caption{问题四 Q4-2（计划线）四个指定日：指定时段计划购电量（kWh）与全天汇总}
\begin{tabular}{lrrrrrrrr}
\toprule
 & \multicolumn{6}{c}{指定时段计划购电量} & 全天购电量 & 全天购电费 \\
\cmidrule(lr){2-7}
日期 & 10 & 12 & 14 & 16 & 18 & 20 & （kWh） & （元） \\
\midrule
03-20 & 0 & 683 & 0 & 474 & 0 & 702 & 58\,953 & 39\,697.50 \\
06-21 & 0 & 0 & 0 & 136 & 412 & 0 & 43\,224 & 21\,074.01 \\
09-23 & 0 & 435 & 0 & 495 & 0 & 18 & 63\,804 & 54\,676.15 \\
12-21 & 0 & 908 & 0 & 840 & 731 & 0 & 93\,979 & 71\,622.82 \\
\bottomrule
\end{tabular}
\label{tab:q4d2}
\end{table}

\begin{table}[H]
\centering\footnotesize\setlength{\tabcolsep}{3pt}
\caption{问题四 Q4-3（调整线）四个指定日：计划与调整购电量（kWh）及全天汇总}
\begin{tabular}{lrrrrrrrrrrrrrr}
\toprule
 & \multicolumn{6}{c}{计划购电量 $g$} & \multicolumn{6}{c}{调整购电量 $a$} & \multicolumn{2}{c}{全天量 $g/a$} \\
\cmidrule(lr){2-7}\cmidrule(lr){8-13}\cmidrule(lr){14-15}
日期 & 10 & 12 & 14 & 16 & 18 & 20 & 10 & 12 & 14 & 16 & 18 & 20 & （kWh） & 购电费（元） \\
\midrule
03-20 & 0 & 543 & 0 & 566 & 0 & 702 & 0 & 543 & 0 & 566 & 8 & 715 & 57\,890 & 40\,213.26 \\
06-21 & 0 & 0 & 0 & 182 & 415 & 0 & 0 & 0 & 0 & 95 & 415 & 0 & 43\,178 & 20\,875.89 \\
09-23 & 0 & 376 & 0 & 492 & 0 & 18 & 0 & 289 & 0 & 492 & 0 & 18 & 65\,806 & 47\,746.03 \\
12-21 & 0 & 995 & 0 & 832 & 726 & 0 & 0 & 934 & 0 & 877 & 690 & 0 & 92\,543 & 71\,641.03 \\
\bottomrule
\end{tabular}
\label{tab:q4d3}
\end{table}

\subsection{模型检验}

\textbf{退化复现（判定性）。}常数价退化检验（附件一价格、常数终端价值、全年 334 天）实现两条交付线对前三问主线的逐位复现：Q4-2 线逐日复现问题二交付结果（13\,745\,816.13 元，334 天全部明细列零误差），Q4-3 线逐日复现问题三交付结果（13\,023\,147.06 元，同口径零误差）。波动价管线与题二、题三共用物理内核与执行层的实现一致性由此闭合。

\textbf{自检。}结算价残差（交付表费用列按实现价逐日重算）为 0；对账与能量守恒残差 $\le10^{-6}$；序关系（末端估值调整口径）零违例；紧急购电上界的序列线性化通过不动点自检；日频超买占比 88.62\%/86.23\% 落入报童区间 $[70\%,90\%]$。当前唯一未通过的判定条款是“放电与弃电不同时发生”（两条线各有 1\,791/1\,232 个时段违反）：源于求解器在等价最优顶点中选择“先放电再弃掉”的路径，对现金费用无影响，已通过给弃电加 $10^{-6}$ 量级微罚收紧，残余时段正在归零，本文如实报告。

% ===================== 九、副模型与对照实验体系 =====================
\section{副模型与对照实验体系}\label{sec:aux}

\subsection{设计原则：一次只改一处}

主模型交付的是一串费用数字，而数字本身说明不了模型好在何处、哪一部分结构真正在起作用。为此本文围绕主模型建立一套副模型体系：每条副模型与主模型\textbf{共用}预报层、情景生成器、执行层实现、跨日 SOC 链与随机种子，\textbf{只改动一处}结构，于是两者的费用差就是对该处结构的干净归因，不掺入实现差异与随机性差异。这一纪律使本文每一个“某机制值多少钱”的结论都可独立复核，也使“改动只让指标好看、总费用却上升”这类伪改进无处藏身。按改动对象，副模型分四类（表 \ref{tab:auxmap}）。

\begin{table}[H]
\centering\small
\caption{副模型体系总览：只改一处，费用差即该处的归因}
\begin{tabular}{lllc}
\toprule
类别 & 副模型 & 只改动的一处 & 回答的问题 \\
\midrule
\multirow{3}{*}{信息类}
 & 完美信息 $\ZPI$ & 0:00 即知当日真值 & 不确定性总共值多少 \\
 & 事后最优 $\Zep$ & 计划锁定后补救全知 & 执行层的极限在哪 \\
 & 读法 D（题四） & 0:00 即知当日实现电价 & 当日价格信息值多少 \\
\midrule
\multirow{3}{*}{策略类}
 & 无电池报童 $Z^{\mathrm{nobat}}$ & 移除储能，只留分位裕度 & 要不要电池 \\
 & 谷充峰放规则 $Z^{\mathrm{rule}}$ & 优化换人工阈值规则 & 要不要优化 \\
 & 题一计划复读 $Z^{\mathrm{Q1}}$ & 不逐日重规划 & 要不要逐日适应 \\
\midrule
\multirow{6}{*}{机制类}
 & 固定 MPC $\ZCE$ & 未来压成单条分位路径 & 未来按分布还是按路径 \\
 & 不改单 E0 & 冻结 $a\equiv g$ & 修订权值多少（附加问） \\
 & 不建模价格 PR0 & 价格预报换典型日 & 电价要不要建模 \\
 & 常数终端价值 V0 & 终端价值不逐日标定 & 终端价值要不要逐日 \\
 & 无口径上界 & 紧急购电放开为自由变量 & 口径收紧的合规成本 \\
 & 独立残差抽样 & 价格残差与负荷残差分开抽 & 情景要不要三元联合 \\
\midrule
\multirow{2}{*}{换账单 / 选臂}
 & 换账单 $\Zmis$（题四） & 不重优化、只换结算价 & 贵出的钱里多少是价格 \\
 & $3\times3\times2$ 预报网格 & 预报候选组合 & 预报层该选哪一臂 \\
\bottomrule
\end{tabular}
\label{tab:auxmap}
\end{table}

\subsection{信息类：费用的上下界与信息阶梯}

同一条真实路径上三个口径满足 $\ZPI\le\Zep\le\Zon$：右不等号因为给定同一计划 $g$ 时事后最优补救的可行域包含任何在线策略的执行结果，左不等号因为完美信息连 $g$ 都在真值下优化。任一处倒序即为实现错误，这是本文最基本的判定性检验（全年零违例）。据此得 $\mathrm{EVPI}=\Zon-\ZPI$，并可按“计划选错”与“补救看不到未来”二分解：问题二 $12\,223\,093\le13\,463\,631\le13\,745\,816$，$\mathrm{EVPI}=152.3$ 万元，日前与实时之比 $81.5:18.5$；问题三 $12\,243\,667\le12\,983\,832\le13\,023\,147$，$\mathrm{EVPI}=77.9$ 万元，实时支仅 $3.9$ 万元。

其中\textbf{实时信息损失}（$\Zon-\Zep$）不含完美信息项，跨问可直接比较：由问题二的 28.2 万元降至问题三的 3.9 万元，说明修订权把原来丢在执行层的钱提前在承诺层挣回来了，问题三剩余的改进空间集中在日前而非执行。$\mathrm{EVPI}$ 本身则不作跨问定量比较——本文的逐日 $\ZPI$ 挂各臂自身在线轨迹的日初储电量，含臂间 SOC 锚差异，仅在同一问内部作二分解使用。

问题四多一级：把“完全忽略价格波动”（换账单 $\Zmis$）与“0:00 即知当日价”（读法 D）补进去，得四级信息阶梯 $\ZPI\le\ZD\le\Zon\le\Zmis$，依次回答全知、只知价、知价格规律（主交付）、忽略波动，两条线实测均成立：计划线 $12\,707\,701\le14\,402\,278\le14\,447\,040\le14\,519\,354$，调整线 $12\,741\,992\le13\,580\,665\le13\,671\,435\le13\,744\,783$。两点口径声明：$\ZD$ 只是\textbf{同期合计}上界而非逐日上界（334 天中有 134 天读法 D 反而更贵——它抽电更狠，独立 SOC 链把成本推到后面几天），故只在全年合计意义下解读；问题四的完美信息记 $Z^{\mathrm{PI|arm}}$，挂本臂日初 SOC，跨臂不可比。

\subsection{策略类：三级笨办法基准}

三条副模型把优化逐级换成人工可执行的做法，共用同一预报、同一执行层与同一结算（数字见表 \ref{tab:q2ladder}）。\textbf{要不要电池}：移除储能后按报童分位留裕度，最优臂仍贵 $24.4\%$，即储能的年化价值。\textbf{要不要优化}：同样带电池，“低于 35\% 价格分位充、高于 65\% 分位放”的阈值规则贵 $16.5\%$。\textbf{要不要逐日适应}：把问题一的典型日计划当全年计划逐日复读，贵 $80.5\%$，紧急购电占其费用的 $52.73\%$——近一半的钱花在 5 倍罚则上，说明“计划必须随日更新”不是锦上添花。

无电池臂的分位扫描另有一重作用：它是报童临界分位公式的独立交叉验证（表 \ref{tab:nobat}）。缺口频率随 $\alpha$ 单调下降（$13.70\%\to2.53\%$），费用却呈 U 形，最低点落在 $\alpha=0.85$，而理论区间为 $\alpha^\star=(\kappa-1)/(\kappa-\theta)\in[0.80,0.96]$（过量回收率 $\theta$ 自 0 变到 $\eta^2$ 的两端），实测最低点正在区间内。该曲线同时是\textbf{过度对冲的反例}：把覆盖水平推过最优点，缺口率确实按 $\kappa$ 的比例塌缩，但边际收益远低于边际成本——这比单报一个最优值更能说明裕度不是越厚越好。

\begin{table}[H]
\centering\small
\caption{无电池报童副模型的分位扫描（334 个交付日）：费用 U 形，最低点落入理论区间}
\begin{tabular}{crrr}
\toprule
分位 $\alpha$ & $Z^{\mathrm{nobat}}$（元） & 相对交付 & 紧急购电费占比 \\
\midrule
0.70  & 17\,399\,386.60 & $+26.6\%$ & $13.70\%$ \\
0.80  & 17\,154\,844.54 & $+24.8\%$ & $9.55\%$ \\
\textbf{0.85} & \textbf{17\,096\,785.82} & $\bm{+24.4\%}$ & $7.48\%$ \\
0.90  & 17\,138\,896.24 & $+24.7\%$ & $5.81\%$ \\
0.955 & 17\,292\,411.80 & $+25.8\%$ & $3.62\%$ \\
1.00  & 17\,565\,835.84 & $+27.8\%$ & $2.53\%$ \\
\bottomrule
\end{tabular}
\label{tab:nobat}
\end{table}

\subsection{机制类之一：未来该按分布还是按路径（固定 MPC）}

固定 MPC 即把执行层剩余时域的情景集压成单条分位参考路径、其余逐字不变的副模型。它与主模型只差“未来是一个分布还是一条路径”，因此费用差是干净的机制归因（表 \ref{tab:aux_ce}）。机制解释见 \ref{sec:q2} 节：分位压缩等价于令缺口发生概率 $q\equiv1$，把“可能缺口”当“必然缺口”给储能定影子价格，于是拒绝放电、改在谷时以 5 倍价囤电。

把固定 MPC 当作交付策略完整跑全年（早期口径，14\,069\,441 元），三条病征同时出现且在换回随机 MPC 后同时收敛：谷价段紧急购电 $51\%$、缺口频率 $11.19\%$ 远超设计值、$86.5\%$ 的紧急电量发生在电池有电时（细节见 \ref{sec:q2} 节检验段）。跨四问还有一条结构性规律：执行层增益随上游信息丰富\textbf{单调收缩}（50.9 万 $\to$ 8.3 万；32.4 万 $\to$ 3.7 万）——上游越有机会把计划定对，留给执行层补救的量越少，与问题三实时信息损失只有 3.9 万元是同一件事的两种度量。

\begin{table}[H]
\centering\small
\caption{固定 MPC 副模型：同计划、同日初 SOC 的配对对照（元）}
\begin{tabular}{lrrr}
\toprule
问题 & 主模型 $\Zon$ & 固定 MPC $\ZCE$ & 执行层增益 \\
\midrule
问题二 & 13\,745\,816.13 & 14\,254\,554.72 & 508\,738.58 \\
问题三 & 13\,023\,147.06 & 13\,106\,182.26 & 83\,035.20 \\
问题四（计划线） & 14\,447\,039.68 & 14\,771\,011.74 & 323\,972.06 \\
问题四（调整线） & 13\,671\,435.06 & 13\,708\,593.11 & 37\,158.06 \\
\bottomrule
\end{tabular}
\label{tab:aux_ce}
\end{table}

\subsection{机制类之二：修订权（不改单副模型）}

不改单副模型冻结 $a\equiv g$、其余逐字不变，是附加问的基准。该基准的合格性有两个必须交代的细节，否则修订权的价值会被系统性夸大：其一，它的计划 $g$ 由\textbf{非预期性}计划器给出，而不是“先用假定后面能改的规划器解出计划、再把它冻住”——后者是稻草人，本文实测该错误基准会把费用虚高 232.7 万元；其二，冻结须可核验，该臂的上调与下调电量合计为 0、偏差费为 0、调整行使率 $0.0\%$，冻结确实生效。

在此基准上逐级开放修订权，两种价格环境下的阶梯如下（常数价即问题三，见 \ref{sec:q3} 节）：常数价 $13\,633\,399\to13\,345\,217\to13\,176\,494\to13\,023\,147$，合计 $610\,252$ 元（$4.48\%$），6:00 档占 $47\%$；波动价（调整线）$14\,323\,490\to14\,006\,509\to\cdots\to13\,671\,435$，合计 $652\,055$ 元（$4.55\%$），6:00 档 $316\,981$ 元占 $49\%$，12:00 与 18:00 两档合计 $335\,074$ 元（该两档的单档回测未完成，故不细分）。读法：修订权的价值在常数价（$4.48\%$）与波动价（$4.55\%$）下几乎一致，且 6:00 档在两种价格环境下都贡献近半——这说明修订权是\textbf{机制本身}的价值，不依赖价格环境，因而附加问的结论（需要引入 6/12/18 时的预报，且越早的修订越值钱）在问题四下同样成立。

\subsection{机制类之三：电价要不要建模、终端价值要不要逐日}

这两处改动必须一起讨论，因为它们在实现上耦合：不建模价格的副模型没有次日价格预报可用，其终端价值会自动退化为常数，于是“$\text{PR0}$ 相对主交付”这一个差额里同时装了两处改动，直接归因给价格预报会高估一倍。本文因此插入只退终端价值的中间臂 V0，把两笔账严格分开（表 \ref{tab:aux_price}）：两条线的分项之和均精确闭合到总差额，确认分解无残项。

\begin{table}[H]
\centering\small
\caption{价格建模与终端价值逐日化的两步分解（334 个交付日，元）}
\begin{tabular}{lrr}
\toprule
臂 / 分项 & Q4-2（计划线） & Q4-3（调整线） \\
\midrule
主交付（价格预报 PR1 $+$ 终端价值逐日 V2） & 14\,447\,039.68 & 13\,671\,435.06 \\
只退终端价值为常数 0.472（V0） & 14\,466\,703.19 & 13\,690\,646.94 \\
再退价格预报为附件一典型价（PR0） & 14\,485\,987.96 & 13\,707\,774.06 \\
\midrule
终端价值逐日化的价值 & 19\,663.51 & 19\,211.88 \\
纯价格点预报的价值 & 19\,284.77 & 17\,127.12 \\
\quad 合计（$=$ PR0 $-$ 主交付，闭合） & 38\,948.28 & 36\,339.00 \\
\bottomrule
\end{tabular}
\label{tab:aux_price}
\end{table}

终端价值逐日化另有一项\textbf{比费用差更硬的结构证据}。V0 臂的日末水位按次日星期分组为 6\,223 / 8\,972 / 5\,840 / 6\,673 kWh（次日为周五 / 周六 / 周一 / 其余），无任何单调台阶；主交付臂则为 2\,594 / 6\,480 / 6\,309 / 6\,974 kWh，次日为周五时显著最低。也就是说，“周四主动清仓、迎接低价周五”这一行为只在逐日标定下出现，常数终端价值在结构上不可能产生它。逐日化的依据因此不止是 1.97 万元的费用差，而是一条被数据证实的行为预言。

\subsection{机制类之四：紧急购电的口径上界}

放开上界是唯一一条\textbf{会让费用变便宜}的副模型：问题二便宜 23\,338 元（$0.17\%$），问题三在同代码仅切换该开关的受控对照中便宜 $0.16\%$。本文仍取带上界的口径，因为便宜的那部分钱来自题面不允许的行为（囤货与盈余时段买入的机理见 \ref{sec:q2cap} 节）。这条副模型给出一条判据：\textbf{当“便宜”来自可行域给多了，便宜不是优势}。

\subsection{机制类之五：情景抽样，三元联合还是独立}

第 \ref{sec:q4} 节的性质 2 断言：价格不确定性的影响集中在协方差通道，独立抽样会低估“高缺口 $\times$ 高价”组合的概率。该论断现有一条直接的副模型证据：把价格残差与负荷、光伏残差分开独立抽样（其余逐字不变），计划线全年费用 14\,454\,917.33 元，比三元联合贵 7\,877.65 元。差额几乎全部落在紧急购电上——独立抽样下紧急购电费 981\,927 元（占 $6.79\%$），联合抽样为 955\,822 元（占 $6.62\%$）：低估价格与缺口的正相关，模型就低估最坏组合的发生概率、少留裕度，最终按 5 倍价把欠的账补回来。这是“情景必须联合抽样”从一行等式到一个实测差额的闭环。

\subsection{换账单副模型：价格效应与策略适应的分离}

换账单副模型不重新优化，只把前三问的交付轨迹搬到附件四实现价下重开账单，得到中间量 $\Zmis$，从而把“题四比题二贵”拆成价格环境与策略适应两格（表 \ref{tab:q4mis}）。原价下重开账单可精确复现原交付费用（闭合残差 $10^{-4}$ 元量级），保证该分解不是记账口径的产物。

\begin{table}[H]
\centering\small
\caption{换账单三格分解（元）：$\Zmis$ 为交付轨迹按实现价重开账单、不重新优化}
\begin{tabular}{lrrrr}
\toprule
线 & 原账单 $Z^{Q}$ & 换价 $\Zmis$ & 重优化 $Z^{Q4}$ & 价格效应 / 策略适应 \\
\midrule
Q4-2 & 13\,745\,816.13 & 14\,519\,353.69 & 14\,447\,039.68 & $+773\,537.55$ / $-72\,314.01$ \\
Q4-3 & 13\,023\,147.06 & 13\,744\,782.86 & 13\,671\,435.06 & $+721\,635.80$ / $-73\,347.80$ \\
\bottomrule
\end{tabular}
\label{tab:q4mis}
\end{table}

两条线的价格效应都在 $+5.5\%$ 至 $+5.6\%$、策略适应都在 $-0.5\%$ 左右。前者是价格环境的必然代价（不可通过建模消除），后者才是“在波动价下重新优化”真正挣到的钱，而它的量级（$0.50\%$）与上节价格消融给出的 $0.27\%$ 相容——差额来自三元联合情景等其余波动价适应。这同时给“电价要不要建模”一个上限意义的回答：可挣的空间本就只有半个百分点，本文的价格预报已拿到其中过半。

\subsection{选臂副模型：$3\times3\times2$ 预报消融网格}

预报层的三个候选轴（负载 3 种 $\times$ 光伏 3 种 $\times$ 残差窗 2 种）共 18 组，全部在同一执行层与结算层下跑完全年，以 $\Zon$ 定夺而不以预报误差定夺（理由见附录 \ref{sec:appBwin}：预报技能提升不等于省钱）。结果见表 \ref{tab:aux_grid}。

\begin{table}[H]
\centering\small
\caption{预报层 18 组消融网格（334 个交付日，节选首末与交付臂）}
\begin{tabular}{clrr}
\toprule
排名 & 组合 & $\Zon$（元） & 相对最优臂 \\
\midrule
1 & \textbf{L2-P1-R2} & \textbf{13\,659\,949.10} & — \\
2 & L2-P2-R2 & 13\,673\,864.61 & $+0.10\%$ \\
3 & L1-P1-R2（本文交付臂） & 13\,745\,816.13 & $+0.63\%$ \\
9 & L3-P1-R2 & 14\,081\,979.94 & $+3.09\%$ \\
18 & L3-P3-R1 & 14\,483\,654.66 & $+6.03\%$ \\
\bottomrule
\end{tabular}
\label{tab:aux_grid}
\end{table}

三条轴的内部序在所有配对中完全一致——负载轴 L2 $>$ L1 $>$ L3（6/6 配对）、光伏轴 P1 $>$ P2 $>$ P3、残差轴 R2 $>$ R1（9/9 配对），最优臂正是三轴的单调顶点，说明这是结构信号而非抽样噪声。机制上三条序各有解释：L2（同星期几指数加权）压过 L1（近 4 个同星期几均值），说明“样本更多”与“更贴近期”之间的连续插值优于硬截断；L3（全历史对数线性回归配星期效用）垫底，证否了它隐含的“星期效用全年恒定”假设，周节律随季节漂移；R2（近 12 日残差全取）9/9 压过 R1（28 日放回抽 25），与负载轴同向，即近期残差比大样本更重要。最优臂的判定性自检（紧急购电归因四条、上界不动点收敛）全部通过；交付臂在网格中复跑与交付结果逐位一致（13\,745\,816.13 元），确认网格与交付出自同一实现。

如实报告：本文\textbf{未}把交付臂换成最优臂。问题三、问题四的预报层与全部对照基准均在 L1-P1-R2 上标定，换臂需要整条链重跑并重算全部副模型；因此 $0.63\%$（85\,867 元）是一处已量化、已定位的可改进空间，而非未知风险。

\subsection{小结：副模型体系给出的四条结论}

\begin{enumerate}[leftmargin=2em]
  \item \textbf{三件事各值多少。}电池值 $24.4\%$、优化值 $16.5\%$、逐日适应值 $80.5\%$——三条策略类副模型把“为什么需要这个模型”量化成了三个数，其中逐日适应的价值最大。
  \item \textbf{未来必须以分布进入模型，且上游越准，执行层越不重要。}固定 MPC 的配对对照在四问全部更贵，且执行层增益随上游信息丰富而单调收缩（50.9 万 $\to$ 8.3 万 $\to$ 3.7 万）。
  \item \textbf{修订权是机制价值，与价格环境无关。}常数价 $4.48\%$ 与波动价 $4.55\%$ 几乎相等，6:00 档在两种环境下都贡献近半。
  \item \textbf{电价的价值分两层且量级相当。}可从历史学到的一层（价格预报加终端价值逐日化）值 $0.27\%$，只有当天才知道的一层（读法 D）值 $0.31\%$，二者接近说明自建预报已接近“能从历史学到的极限”，也说明就算完全知道当日电价，也只能再省三个千分点。
\end{enumerate}

% ===================== 十、模型评价与改进 =====================
\section{模型评价与改进}

\subsection{优点}

\begin{enumerate}[leftmargin=2em]
  \item \textbf{四问共骨架，特例关系可检验。}物理可行域四问共用，后一问的可行域包含前一问，使“费用不得超过前一问”“单情景退化逐位复现”成为判定实现正确性的硬检验。本文开发过程中正是由该检验发现并修复了问题三日前规划程序的一处变量布局缺陷（修复使全年费用下降 93.0 万元），相关复现检验已固化为版本库中的自动化测试；
  \item \textbf{诚实回测保证结论可信。}所有预报与情景只用决策时刻之前的数据，附件二/四按日解锁；十余条对照基准共用同一预报层、执行层与随机种子，每条只改动一处结构，因此“某机制值多少钱”的每个结论都是干净归因而非混合效应（副模型体系见 \ref{sec:aux} 节）；
  \item \textbf{理论量贯穿始终。}套利门槛 $1/\eta^2$、报童临界分位（缺口频率 4.30\% 对理论 4.5\%，日频超买 88.6\% 对 80\%）、隔夜水位、星期台阶等先验预言均在回测中量化对照，模型的每个结构性结论都有机制解释；
  \item \textbf{口径与代价公开。}紧急购电上界的合规成本（题二 $+0.17\%$、题三 $+0.16\%$）、弃电拆项中“报童保险费不是浪费”的解读、未通过的判定条款，均如实在文中报告。
\end{enumerate}

\subsection{缺点}

\begin{enumerate}[leftmargin=2em]
  \item 日前规划的第二阶段在情景内假定“届时知道哪个情景成真”（等待型决策），是乐观近似，其精确化需要情景树结构，本文仅以“非预期性计划 + 修订权后置”规避并量化了该代价（15.1 万元）；
  \item 执行层的情景集沿用 0:00 版本、不做日内更新（诚实回测允许用当日已实现值校正），保守成分集中在真缺口一侧；
  \item “放电与弃电互补”条款在题四尚有 1\,791/1\,232 个时段未归零（等价最优顶点问题，不改现金费用），是当前唯一未过的自检；
  \item 题二的 18 组预报网格显示最优臂为 L2-P1-R2，较交付臂便宜 $0.63\%$（85\,867 元），本文未换交付臂——问题三、问题四的预报层与全部副模型均在交付臂上标定，换臂须整链重跑；该差距已量化、已定位（\ref{sec:aux} 节）；
  \item 题三的十二组预报消融网格在修正后计划器下重跑未完成，其预报层配置按机制理由并参照题二网格的同轴结论选定；题四修订阶梯的 12:00 与 18:00 两档仅有合计数、单档回测未细分；
  \item 调整行为中下调电量多于上调（0.53 倍），“宁存不退”的理论方向在交付解上未获验证，说明日前计划对日内预报下修的敏感性仍可改进。
\end{enumerate}

\subsection{改进方向}

将执行层情景集按当日已实现残差每 6--12 个时段重建（预期同时降低在线费用与实时信息损失）；日前规划改为按修订时刻分组的情景树共享变量，消除等待型近似；在确认换臂不影响题三、题四继承链后，评估是否把预报层交付臂换成网格最优的 L2-P1-R2；补齐题三十二组消融、题四修订阶梯的单档细分，以及放电—弃电微罚的全时段收口。

\subsection{检验结论汇总}

能量守恒与对账恒等式残差四问全部 $\le10^{-6}$；序关系（题四用末端估值调整口径）全年零违例；退化复现（题三至题二、题四至题二与题三、题二至题一）逐位通过；报童频率检验两个口径（时段缺口频率 4.30\%、日频超买 88.6\%/86.2\%）均落入理论区间；末端估值敏感性 $+0.79\%\sim+1.43\%$；充放互斥、紧急购电不同时充电等结构性质逐时段核验为零违例。

% ===================== AI 声明（固定句式，位置在参考文献之前） =====================
\section*{AI 工具使用声明}
本参赛队在竞赛过程中使用了 AI 工具，主要用于建模思路讨论、代码实现与调试、论文初稿撰写，详细使用情况见支撑材料。

% ===================== 参考文献 =====================
\begin{thebibliography}{9}\small
\bibitem{jiang2018} 姜启源, 谢金星, 叶俊. 数学模型[M]. 5版. 北京: 高等教育出版社, 2018.
\bibitem{sikou} 司守奎, 孙兆亮. 数学建模算法与应用[M]. 3版. 北京: 国防工业出版社, 2021.
\bibitem{birge} Birge J R, Louveaux F V. Introduction to Stochastic Programming[M]. 2nd ed. New York: Springer, 2011.
\bibitem{rawlings} Rawlings J B, Mayne D Q, Diehl M. Model Predictive Control: Theory, Computation, and Design[M]. 2nd ed. Madison: Nob Hill Publishing, 2017.
\bibitem{qin2011} Qin Y, Wang R, Vakharia A J, et al. The newsvendor problem: Review and directions for future research[J]. European Journal of Operational Research, 2011, 213(2): 361--374.
\bibitem{bates1969} Bates J M, Granger C W J. The combination of forecasts[J]. Journal of the Operational Research Society, 1969, 20(4): 451--468.
\bibitem{sioshansi} Sioshansi R, Denholm P, Jenkin T, et al. Estimating the value of electricity storage in PJM: Arbitrage and some welfare effects[J]. Energy Economics, 2009, 31(2): 269--277.
\bibitem{denholm2007} Denholm P, Margolis R M. Evaluating the limits of solar photovoltaics in traditional electric power systems[J]. Energy Policy, 2007, 35(5): 2852--2861.
\bibitem{huangfu2018} Huangfu Q, Hall J A J. Parallelizing the dual revised simplex method[J]. Mathematical Programming Computation, 2018, 10(1): 119--137.
\end{thebibliography}

% ===================== 附录 =====================
\newpage
\appendix

\section*{附录 A\quad 支撑材料与图件}
附录 B 给出终端价值与预报窗口的标定实验，附录 C 收录核心建模源代码。

\begin{table}[H]
\centering\small
\caption{支撑材料清单}
\begin{tabular}{ll}
\toprule
路径 & 内容 \\
\midrule
\texttt{src/nmmc/} & 全部求解源代码（数据 / 模型 / 回测 / 报表 / 命令行入口） \\
\texttt{experiments/} & 诊断与对照实验脚本（题二、题三、缺口汇总） \\
\texttt{tests/} & 判定性复现检验（题三、题四自动化测试） \\
\texttt{results/} & 各问交付结果与明细（result1--result4 等）；\texttt{results/paper/} 为论文图件 \\
\texttt{docs/} & 建模文档（变量定义、修正报告、数据核验、图件清单） \\
\bottomrule
\end{tabular}
\end{table}

复现命令（仓库根目录，Python 3.12）：问题一至四分别为 \texttt{python -m nmmc.cli.prob1\,\textbf{--}\,output-dir results/q1} 及 prob2/prob3/prob4 同构命令；判定性测试 \texttt{python -m unittest tests.q3.test\_c3\_gates tests.q4.test\_c4\_gates}；论文图重绘 \texttt{python -m nmmc.cli.paper}。

附录图为题三全年修订指纹热图（图 \ref{fig:qA2}）：红色为上调、蓝色为下调，0--6 时块法定空白（无修订权），调整集中于午后可修订块。

\begin{figure}[H]
\centering
\includegraphics[width=0.80\textwidth]{figA2_题3修订指纹-红上调蓝下调-B0法定空白.png}
\caption{题三全年修订指纹（334 天 $\times$ 144 时段）：红 $=$ 上调，蓝 $=$ 下调，夜间块法定空白}
\label{fig:qA2}
\end{figure}

\section*{附录 B\quad 外生参数与超参数标定}
\setcounter{section}{2}
\setcounter{subsection}{0}
\renewcommand{\thesubsection}{B.\arabic{subsection}}

\subsection{全模型外生量清单}

本文刻意把“需要人来定的数”压到最少。表 \ref{tab:hyper} 列出模型中全部外生量及其来历：绝大多数直接由题面与附录一给定，真正需要标定的只有终端价值与两个预报窗口，且三者都做了专门的验证实验（\ref{sec:appBnu}、\ref{sec:appBwin}）。

\begin{table}[H]
\centering\small
\caption{外生量清单：取值、来历与验证状态}
\begin{tabular}{llll}
\toprule
量 & 取值 & 来历 & 验证 \\
\midrule
$\Delta$ & $1/6$ h & 附件数据粒度 & 不适用 \\
$\eta,\ C^{\max}{=}D^{\max},\ [\underline S,\overline S]$ & 0.90，833.33 kWh，$[1200,10800]$ & 附录一给定 & 不适用 \\
$\kappa,\ \beta^-,\ \beta^+$ & 5，0.5，1.5 & 题面罚则 & 不适用 \\
$S_0^{(1)}$ & 6\,000 kWh & 附录一标定（题二起为全年链锚） & 题一扫描，图 \ref{fig:q1s0} \\
\midrule
$\nuend$ & 0.472 元/kWh（题四逐日） & 题一循环约束对偶，非拟合 & 三重验证（\ref{sec:appBnu}） \\
$W_L$ & 4（同星期几步） & 负载周节律窗 & 去超参实验（\ref{sec:appBwin}） \\
$W_R$ & 7 日 & 光伏持续性窗 & 去超参实验（\ref{sec:appBwin}） \\
$N_s$ & 12 & 残差窗全取，等概率 & 网格 R 轴 9/9 胜 R1 \\
\midrule
$\varepsilon$ & $10^{-6}$ & 弃电微罚与数值容差 & 量级远低于费用精度 \\
负载缩放截断 & $[0.7,1.3]$ & 题三日内比例更新的保护带 & 防止异常日把缩放拉飞 \\
随机种子 & 0 & 全流程固定 & 结果可逐位复现 \\
$\alpha^{\mathrm{rt}}$ & 0.65（\textbf{仅}固定 MPC 副模型） & 主模型中已被消去 & 冻结不扫参 \\
\bottomrule
\end{tabular}
\label{tab:hyper}
\end{table}

关于 $S_0$：问题一把它决策化（只要求昼夜循环且落在运行区间内），最优值 8\,550 kWh 由模型给出；若固定为附录一标定的 6\,000 kWh，费用 $Z^{6000}=35\,126.95$ 元，仅贵 8.35 元（$+0.024\%$），全区间扫描的费用带宽也只有 $0.24\%$（图 \ref{fig:q1s0}）。因此“把 $S_0$ 交给模型”既不是取巧也不敏感。问题二起 $S_0^{(1)}=6\,000$ 恢复为给定状态，作全年 SOC 链的锚。

\subsection{终端价值 $\nuend=0.472$ 的来历与三重验证}\label{sec:appBnu}

\textbf{来历（推导而非调参）。}$\nuend$ 回答“今晚在池里留 1 kWh 值多少钱”。问题一昼夜循环约束的对偶变量恰是这个量，数值 0.4720 元/kWh；其经济学可行区间为 $[\,p^{\mathrm v}/\eta,\ \eta p^{\mathrm p}\,]=[0.4126,1.2557]$（下端为谷价重置成本、上端为峰价变现收益），区间非空当且仅当峰谷比不低于套利门槛 $1/\eta^2$，与问题一的性质 2 对接。

\textbf{验证一：局部重扫。}在交付口径下取 20 天窗口逐点重跑（表 \ref{tab:nu}），费用随 $\nuend$ 单调上升，锚点最优。

\begin{table}[H]
\centering\small
\caption{$\nuend$ 局部重扫（交付口径，20 天窗口）}
\begin{tabular}{crrr}
\toprule
$\nuend$ & 窗口费用（元） & 相对锚点 & 紧急费占比 \\
\midrule
\textbf{0.472} & \textbf{615\,622} & — & $9.40\%$ \\
0.60 & 620\,466 & $+0.79\%$ & $9.83\%$ \\
0.80 & 622\,507 & $+1.12\%$ & $10.10\%$ \\
1.00 & 624\,405 & $+1.43\%$ & $10.30\%$ \\
\bottomrule
\end{tabular}
\label{tab:nu}
\end{table}

\textbf{验证二：全年地标扫描。}在无口径上界的早期全年口径下，取理论区间下端、锚点与两个偏高地标各跑全年：0.4126 得 13\,736\,047 元、\textbf{0.472 得 13\,722\,478 元}、0.616 得 13\,773\,077 元、0.7662 得 13\,803\,862 元。最小值恰在锚点、两侧均变贵，四点极差 $0.742\%$。交付口径下的局部重扫（验证一）方向相同，故该结论对上界不敏感。

\textbf{验证三：可达性与逐日化。}把全年每一天的次日曲线各解一次循环线性规划、取对偶，得到逐日序列：均值 0.4706、标准差 0.0049、最大值 0.4720，365/365 天落在理论区间内。这说明常数价下 0.472 同时是理论值与经验上界，偏高的两个地标（$\ge0.616$）经验上不可达；把常数锚整体换成该逐日序列，全年费用只变化 $0.149\%$，即模型对这唯一的人工参数不敏感。问题四把它升级为逐日标定（以次日预报价解循环线性规划取对偶），至此全模型不含任何拟合常数，首日在附件一价格下复现 0.472 作为标定自检。

\subsection{预报窗口 $W_L=4$ / $W_R=7$ 的合理性}\label{sec:appBwin}

这两个数是全模型仅存的“经验窗口”，因此本文专门做了一组去超参实验来检验它们：把固定窗口换成\textbf{滚动伪样本外选窗}——候选 $W_L\in\{1,\dots,10\}$（同星期几步）、$W_R\in\{1,\dots,14\}$（日），在只含过去日的公共评估集上算平均绝对误差，按 1-SE 规则在带内取最小窗，评估日不足 21 天时退回先验值。四条自适应臂与固定窗的对照见表 \ref{tab:win}（费用为无口径上界的早期全年口径，基线 13\,722\,478 元；此处只用相对量）。

\begin{table}[H]
\centering\small
\caption{窗口去超参实验：预报误差（kWh/时段）与全年费用}
\begin{tabular}{lrrrr}
\toprule
臂 & 负载 MAE & 光伏 MAE & 净负荷 MAE & $\Zon$ 相对 \\
\midrule
固定 4 / 7（本文采用） & 31.480 & 25.567 & 47.484 & — \\
仅负载自适应 $W_L$ & 29.164 $(-7.36\%)$ & 25.567 & 45.773 $(-3.60\%)$ & $+0.60\%$ \\
仅光伏自适应 $W_R$ & 31.480 & 25.774 $(+0.81\%)$ & 47.845 $(+0.76\%)$ & $-0.26\%$ \\
双自适应 & 29.164 & 25.774 & 46.090 $(-2.94\%)$ & $+0.54\%$ \\
零参数（全历史回归 $\times$ 全历史形态） & 41.348 $(+31.4\%)$ & 60.333 $(+136\%)$ & 85.459 $(+80.0\%)$ & 不具可比性 \\
\bottomrule
\end{tabular}
\label{tab:win}
\end{table}

三条结论：

\begin{enumerate}[leftmargin=2em]
  \item \textbf{选窗机制确实在工作}，不是摆设：报告期内负载有 274/334 天、光伏有 330/334 天选到了先验值以外的窗口；误差曲面与 1-SE 选择经双实现交叉对拍（曲面最大偏差 $4.5\times10^{-13}$，逐日选择完全一致）。
  \item \textbf{预报技能提升不等于省钱}，这是本实验最有价值的发现。负载侧平均绝对误差下降 $7.36\%$，全年费用反而上升 $0.60\%$；光伏侧误差上升 $0.81\%$，费用反而下降 $0.26\%$。原因是技能改善未必落在决定计划购电量的关键时段上——真正起作用的是误差落点与经济结构（峰谷位置、储能容量与功率的紧约束时段）的交互，而不是误差的总量。这直接决定了本文的选臂判据：预报层一律以全年费用 $\Zon$ 定夺，不以误差指标定夺（第 \ref{sec:aux} 节的 18 组网格即按此执行）。
  \item \textbf{彻底零参数不可行}：把窗口完全取消、改用全历史回归配全历史形态，净负荷误差恶化 $80\%$。
\end{enumerate}

因此本文保留固定 4 / 7，并把上述现象如实公开：这不是“没试过自适应”，而是自适应在费用口径上不占优。同款检验在问题三的融合层复核一次：固定 $W_R=7$ 对逐日选窗，四个修订阶段的光伏误差分别劣化 $1.03\%$、$0.80\%$、$0.26\%$、$0.24\%$，全阶段平均 $0.76\%$，固定窗同样不劣。

\section*{附录 C\quad 核心建模源代码}
以下收录四问共用的 LP 骨架与各问规划层装配（完整代码含数据层、回测编排、报表与全部实验脚本，见支撑材料）。

\lstinputlisting[caption={\ttfamily src/nmmc/model/lp.py：四问共用 LP 骨架（物理常数、SOC 递推块、求解封装）}]{../src/nmmc/model/lp.py}
\lstinputlisting[caption={\ttfamily src/nmmc/model/q1.py：问题一典型日线性规划}]{../src/nmmc/model/q1.py}
\lstinputlisting[caption={\ttfamily src/nmmc/model/q2.py：问题二两阶段场景规划与随机 MPC 子问题}]{../src/nmmc/model/q2.py}
\lstinputlisting[caption={\ttfamily src/nmmc/model/q3.py：问题三分块承诺规划}]{../src/nmmc/model/q3.py}
\lstinputlisting[caption={\ttfamily src/nmmc/model/q4.py：问题四终端价值逐日标定与价格折叠}]{../src/nmmc/model/q4.py}

\end{document}
